\documentclass[french]{book}
\usepackage{teach}
\usepackage{verbatim}
\usepackage{amsmath,amsfonts,amssymb}
\usepackage{pgfplots}
\usepackage{mwe}
\usepackage{stmaryrd}
\usepackage{yhmath} 
\usepackage{pstricks,pst-plot,pst-text,pst-tree,pst-eucl}
\newcommand{\bbr}{\mathbb {R}}
\newcommand{\bbc}{\mathbb {C}}
\newcommand{\bbz}{\mathbb {Z}}
\newcommand{\I}{i}
\newcommand{\tm}{\times}

\newcommand{\ofr}[2]{\dfrac{#1}{#2}}
\newcommand{\arc}[1]{\wideparen{#1}}
\newcommand{\ab}[1]{\vert #1 \vert}
\newcommand{\nov}[1]{\Vert \overrightarrow{#1} \Vert}

%\newcommand{\arc}[1]{\wideparen{#1}}
\RequirePackage{graphicx}
\RequirePackage{ifpdf}
 \ifpdf
   \DeclareGraphicsRule{*}{mps}{*}{}
 \fi

\newcommand{\sysd}[2]{%
       \left\{
       \begin{aligned}
          #1\\
          #2\\
       \end{aligned}
       \right.%
       }
  \usepackage{fancyhdr}
  \usepackage{tkz-tab}
  \usepackage{amsmath}
  \usepackage{amssymb}
  \usepackage{amsfonts}
  \usepackage{amsthm}
  \usepackage{mathrsfs}

  \usepackage{array}
  \usepackage{multicol}
  \setlength{\multicolsep}{3pt}

  \usepackage{graphicx}
  \graphicspath{{Figures/}}
  \usepackage{pstricks,pst-plot,pst-text,pst-tree,pst-eucl}
  \usepackage[all]{xy}
  \newcommand{\bsyc}{\left\{\begin{array}{rcl}} % syst�e serr�sur le �al
\newcommand{\esyc}{\end{array}\right.}


  \usepackage{fancybox}
  \usepackage{color}

%  \usepackage[frenchb]{babel}
  \usepackage{hyperref}
\usepackage{textcomp}
\usepackage{mathtools}
\usepackage{subfig}
\pgfplotsset{compat=newest}
\usepackage[all]{xy}
%\usepackage{geometry}
%\geometry{top=1cm, bottom=1cm, left=2cm, right=2cm}
\usepackage[utf8]{inputenc}
\usepackage[french]{babel}
\redefinecolor{thm}{title}{bgcolor}{alizarin}
\redefinecolor{prop}{title}{bgcolor}{meatbrown}
\redefinecolor{defn}{title}{bgcolor}{olivine}
\definecolor{alizarin}{rgb}{0.82, 0.1, 0.26}
\definecolor{meatbrown}{rgb}{0.9, 0.72, 0.23}
\definecolor{olivine}{rgb}{0.6, 0.73, 0.45}
\usepackage{mathrsfs}
%%
\newcommand{\R}{\mathbb {R}}
%\newcommand{\C}{\mathds {C}}
\newcommand{\N}{\mathbb {N}}
\newcommand{\Z}{\mathbb {Z}}
\newcommand{\Q}{\mathbb {Q}}
\newcommand{\D}{\mathbb {D}}
\newcommand{\HRule}{\rule{\linewidth}{0.5mm}}
%%
\newcommand{\se}{\geq}
\newcommand{\ie}{\leq}
%\newcommand{\tm}{\times}
%\newcommand{\ell}{l}
%\newcommand{\ell'}{l'}
%%
\newcommand{\barre}[1]{\overline{#1}}
\newcommand{\ds}{\displaystyle}
\newcommand{\limn}{\lim_{n\rightarrow +\infty}}
\newcommand{\limo}[1][x]{\ds\lim_{#1\rightarrow 0}}
\newcommand{\limite}[2][x]{\ds\lim_{#1\rightarrow #2}}
\newcommand{\limpinf}[1][x]{\ds\lim_{#1\rightarrow +\infty}}
\newcommand{\limminf}[1][x]{\ds\lim_{#1\rightarrow -\infty}}
\newcommand{\limiteg}[2][x]{\ds\lim_{#1\xrightarrow{<}#2}}
\newcommand{\limited}[2][x]{\ds\lim_{#1\xrightarrow{>}#2}}
%%
\newcommand{\vect}[1]{\overrightarrow{#1}}
\newcommand{\oij}{(\textit{O},{\overrightarrow{i}},{\overrightarrow{j}})}
%
\newcommand{\co}[2]{C_#2^#1}
\newcommand{\cp}[2]{%
        \begin{pmatrix}
        #1\\
        #2
        \end{pmatrix}%
        }
%
\newcommand{\calb}{\mathscr{B}}
\newcommand{\calc}{\mathscr{C}}
\newcommand{\cald}{\mathscr{D}}
\newcommand{\cale}{\mathscr{E}}
\newcommand{\calf}{\mathscr{F}}
\newcommand{\calp}{\mathscr{P}}
\newcommand{\cals}{\mathscr{S}}
\newcommand{\calt}{\mathscr{T}}
\newcommand{\calr}{\mathscr{R}}


\newcommand{\cali}{\mathscr{I}}
\newcommand{\RR}{\calr }
\newcommand{\CR}{\calc}
\newcommand{\IR}{\cali }

\newcommand{\fr}[2]{\dfrac{#1}{#2}}
\newcommand{\ve}[1]{\overrightarrow{#1}}
\newcommand{\pa}[1]{(#1)}
\newcommand{\ps}[2]{\overrightarrow{#1}.\overrightarrow{#2}}
%\newcommand{\bbr}{\mathbb{R}}
\newcommand{\al}{\alpha}
\newcommand{\la}{\lambda}
\newcommand{\DR}{\cald}
\newcommand{\PR}{\calp}

\newcommand{\nor}[1]{ \Vert #1 \Vert}
\newcommand{\abv}[1]{\vert #1 \vert}
%\newcommand{\coordpp}[3]{\begin{pmatrix}
%
\newcommand{\suite}[1][u]{\left( #1_{n} \right)_{n\in\N}}
\newcommand{\suitep}[1][u]{\left( #1_{n} \right)_{n\in\N^{*}}}
\usetikzlibrary{arrows}
\newcommand{\ssi}{\Longleftrightarrow}
\newcommand{\poly}[1][a]{#1_0+#1_1x+\cdots+#1_nx^n}
\newcommand{\coordpp}[3]{\begin{pmatrix}
#1\\
#2\\
#3
\end{pmatrix}}
\newcommand{\anglevec}[2]{(\overrightarrow{#1};\overrightarrow{#2})}
\newcommand{\un}{\suite}
\newcommand{\eit}{e^{it}}
\newcommand{\eimt}{e^{imt}}
\newcommand{\vn}{\suite[v]}
%\newcommand{\pa}[1]{(#1)}
\newcommand{\cro}[1]{[#1]}
\newcommand{\abs}[1]{\vert #1 \vert}

\newcommand{\Syst}[2]{\left\{\begin{array}{lcccc} #1 \\ #2 \end{array}\right.}
\renewcommand{\arraystretch}{1.2} 
\newcolumntype{C}[1]{>{\centering\arraybackslash}p{#1cm}}
\newcommand{\Rep}{(O;\ve{i};\ve{j})}
\newcommand{\C}[2]{C_{#1}^{#2}}
\newcommand{\dx}{dx}
%\newcommand{\co}[2]{C_#2^#1}
\newcommand{\ud}[1]{\dfrac{#1}{2}}
\newcommand{\cacher}[1]{#1}
\newcommand{\cacheg}[1]{#1}
\newcommand{\cacheb}[1]{#1}
\begin{document}


\begin{titlepage}
  \begin{sffamily}
  \begin{center}
    \textsc{\LARGE Mathématiques expertes}\\[2cm]
    % Title
    { \huge \bfseries Livret de terminal  \\[0.4cm] }
    \HRule \\[2cm]
    \includegraphics[scale=0.17]{../Figures/fill3.1}
    \\[2cm]
    \vfill
    % Bottom of the page
     \large Pierre \textsc{Thiebeaux} \\
    {\large Année scolaire 2019 - 2020}
  \end{center}
  \end{sffamily}
\end{titlepage}

\tableofcontents

\chapter{Les nombres complexes}
\textit{Les  nombres  complexes  portent  bien  leur  nom  !  Ils  interviennent
partout  : en  algèbre, en  analyse, en  géométrie, en  électronique, en
traitement du signal,  en musique, etc. Et en plus,  ils n'ont jamais la
même  apparence  :  tantôt  sous  forme algébrique,  tantôt  sous  forme
trigonométrique, tantôt sous forme  exponentielle, ... Leur succès vient
en fait de  deux propriétés : en travaillant  sur les nombres complexes,
tout polynôme admet un nombre de racines égal à son degré et surtout ils
permettent de calculer  facilement en dimension 2.}


\section{Approche \og moderne\fg{}}

Vous savez compter en dimension 1, c'est  à dire additionner
et multiplier des nombres réels qu'on peut représenter sur la droite des
réels : 

\bigskip

\begin{center}
\includegraphics{../Figures/complexe.1}
\end{center}


Faute  d'outils  plus rigoureux, on  vous a
présenté en classe  de seconde l'ensemble des nombres  réels comme étant
l'ensemble des abscisses des points de la droite orientée ci-dessus. 

Vous  utilisez depuis  l'école primaire  ces nombres  et  les opérations
usuelles qui leur sont  associées, addition et multiplication, sans trop
vous  poser  de  questions.\\
 
 Rappelons  quelques  propriétés \footnote{Ces
  propriétés donnent à $\bbr$ une structure de \textit{corps}, dont voici les propriétés.}
  
\begin{itemize}
\item[$\bullet$] L'addition possède un élément neutre noté 0 : $x+0=0+x=x$.
\item[$\bullet$] La somme de 2 réels est encore un réel.
\item[$\bullet$] Chaque réel $x$ admet un opposé $-x$ vérifiant $x+(-x)=(-x)+x=0$.
\item[$\bullet$]  La multiplication  possède un  élément neutre  noté 1  : $x\times
  1=1\times x=x$ 
\item[$\bullet$] Le produit de deux réels est encore un réel.
\item[$\bullet$] Chaque  réel différent  de 0 admet  un inverse  $x^{-1}$ vérifiant
  $x\times x^{-1}=x^{-1}\times x=1$ 
\item[$\bullet$]   La   multiplication    est   distributive   sur   l'addition   :
  $x\times(y+z)=x\times y + x\times z$. 
\end{itemize}

Tout ceci est bien naturel. Maintenant, on voudrait décoler de l'axe des
réels  et faire  le même  travail en  dimension 2,  c'est-à-dire pouvoir
calculer avec des couples de nombres du style $(x,y)$. Ca veut dire qu'on travaillerai maintenant sur le plan tout entier que l'on  note   $\bbr^2$  cet  ensemble  est  l'ensemble  des coordonnées des points du plan.

La question est: -Est-ce qu'on peut définir une addition
et une multiplication qui  engloberaient et généraliseraient celles vues
dans $\bbr$ et donc nous donnerai un corps plus grand? Une sorte d'"extension de corps". 

Nous allons donc adopter une tactique pour faire cela.  \'A chaque  élément  $(x,y)$ de  $\bbr^2$  nous allons  faire correspondre un nombre qu'on qualifiera de \textit{complexe}. 

L'idée    vient    de   l'observation    \textit{intuitive}\footnote{Les
  $\curvearrowright$  renvoient à une  notion extrêmement  importante et
  rigoureuse  : la notion  d'\textit{isomorphisme}. Deux  ensembles sont
  isomorphes lorsqu'il existe une bijection  entre les deux et que cette
  bijection  <<~conserve~>> les  opérations. Cela  permet  de travailler
  \textit{à  isomorphisme   près}  sur  un  ensemble   compliqué  en  le
  remplaçant   par   un  ensemble   isomorphe   plus   approprié  à   la
  situation.  C'est ce  qui se  passe entre  $\bbr^2$ et  $\bbc$}  : \[\og
(x,y)\curvearrowright  x\cdot(1,0)+y\cdot (0,1)\curvearrowright x\times
1+y\times \sqrt{-1}\curvearrowright x+y\sqrt{-1}\fg{}\]


Nous allons même  donner un nom à ce $\sqrt{-1}$  : appelons-le $i$ pour
qu'il fasse moins peur. Ainsi nous avons les correspondances 

\begin{center}
\begin{tabular}{ccccc}
Le   point   M   &$\longleftrightarrow$    &   Le   couple   $(x,y)$   &
$\longleftrightarrow$ & Le nombre complexe $x+iy$\\ 
$\updownarrow$ & &  $\updownarrow$ & & $\updownarrow$ \\
Le plan $\PR$  & $\longleftrightarrow$ & $\bbr^2$ &$\longleftrightarrow$
& L'ensemble des nombres complexes 
\end{tabular}
\end{center}


Pour se simplifier la vie, nous  allons donner un nom à cet ensemble des
nombres complexes : $\bbc$\\

Et  maintenant   observez  comme  les  calculs   deviennent  faciles  en
prologeant les règles valables sur $\bbr$.

$$(x+iy) + (x'+iy')= x+iy+x'+iy'=(x+x')+i(y+y')$$
$$(x,y)+(x',y')=(x+x',y+y')$$
$$(x+iy)\cdot(x'+iy')=xx'+ixy'+iyx'+i^2yy$$
$$(x+iy)\cdot(x'+iy')=(xx'-yy')+i(xy'+yx')$$


Donc nous  allons pouvoir  calculer en dimension  2 en  généralisant les
règles de  dimension 1. Nous avons  juste ajouté ce nombre  $i$ de carré
$-1$.   En  particulier,  tous   les  nombres  réels  sont  des  nombres
complexes : $\bbr\subset\bbc$.  

Nous allons pouvoir  associer à chacun de ces nombres  réels un point du
plan et  donc associer  des transformations du  plan à des  calculs dans
$\bbc$ : on va résoudre des problèmes de géométrie par le calcul. 
 

\section{Vocabulaire et premières propriétés}




  \begin{thm}[Ensemble $\bbc$]
On  définit un ensemble $\bbc$ 
\begin{itemize}
\item[$\bullet$] muni d'une addition et d'une multiplication qui prolongent celles de $\bbr$
\item[$\bullet$] contenant un nombre $i$ vérifiant $i^2=-1$ 
\item[$\bullet$]  tel que chaque  élément $z$  de $\bbc$  peut s'écrire  de manière
  \textbf{unique}  sous   la  forme   $$z=a+ib  \quad  \text{avec   }  a
  \text{ et } b \text{ des nombres réels}$$ 
\end{itemize}
    
    \end{thm}
  



\subsection{Forme algébrique}

Cette écriture unique est appelée \textbf{forme algébrique} du réel $z$.

Le nombre $a$ est appellé \textbf{partie réelle} de  $z$ et notée $\RR e(z)$

Le nombre $b$ est appellé \textbf{partie imaginaire} de  $z$ et notée $\IR m(z)$


  \begin{rem}
    $\IR m(z)$ est un nombre réel.
    \end{rem}
  


  
    \begin{rem}{\`A quoi sert l'unicité de la forme algébrique ?}
Par exemple,  après maints calculs  savants, vous arrivez au  résultat $
2x+3y-5+i(7x-32y+1)=0$ avec $x$ et $y$  des réels.  Et bien le membre de
gauche  est   une  forme   algébrique  puisque  de   la  forme   réel  +
$i\cdot$réel. Or la forme algébrique de 0 est $0+i\cdot0$.  

Ainsi,  une   équation  complexe   revient  à  deux   équations  réelles
( bienvenue dans la deuxième dimension... ) et donc 

$    2x+3y-5+i(7x-32y+1)=0   \Longleftrightarrow   \begin{cases}
    2x+3y-5=0\\ 7x-32y+1=0\end{cases}$
\end{rem}


\begin{thm}[La formule du binôme]
Soit $z,z'\in \bbc$ alors, 

$$(z+z')^n=\sum_{k=0}^n C^k_nz^kz'^{n-k}$$

\end{thm}

\begin{demo}[Par récurrence sur $\mathbb{N}$]
Initialisation:
$$(z+z')^0=1$$
Hérédité: supposons que $(z+z')^n=\sum_{k=0}^n C^k_nz^kz'^{n-k}$,
$$(z+z')^{n+1}=(z+z')(z+z')^n=(z+z')\sum_{k=0}^n C^k_nz^kz'^{n-k}$$
$$=z\sum_{k=0}^n C^k_nz^kz'^{n-k}+z'\sum_{k=0}^n C^k_nz^kz'^{n-k}$$
$$=\sum_{k=0}^n C^k_nz^{k+1}z'^{n-k}+\sum_{k=0}^n C^k_nz^kz'^{n-(k+1)}$$
$$=\sum_{k=1}^{n+1} C^{k-1}_nz^{k}z'^{n-k-1}+\sum_{k=0}^n C^k_nz^kz'^{n-(k-1)}$$
$$=C^n_n z^{n+1} + \sum_{k=1}^{n} C^{k-1}_nz^{k}z'^{n-k+1} + C^0_nz'^{n+1} + \sum_{k=1}^n C^k_nz^kz'^{n-(k-1)} $$
$$= C^0_nz'^{n+1} + C^n_n z^{n+1} + \sum_{k=1}^{n} (C^{k-1}_n +C^k_n)z^{k}z'^{n-k+1}$$
De plus $C^{k-1}_n +C^k_n=C^k_{n+1}$ d'après la formule de Pascal, d'où
$$=C^0_{n+1}z'^{n+1}+ \sum_{k=1}^{n} C^k_{n+1} z^{k}z'^{n+1-k} + C^{n+1}_{n+1}z^{n+1}$$ 

$$=\sum_{k=0}^{n+1} C^k_{n+1} z^{k}z'^{n+1-k}$$
Par conséquent , notre hypothése de récurence entraine le résultat au rang $n+1$.

Conclusion : donc pour tout $n\in \mathbb{N}$, $$(z+z')^n=\sum_{k=0}^n C^k_nz^kz'^{n-k}$$

\end{demo}
\subsection{Le plan complexe}

Nous avons vu que chaque nombre complexe peut être associé à un point du
plan qu'on munit d'un repère $\pa{O,\ve{e_1},\ve{e_2}}$ 


\begin{figure}
\begin{center}
  \includegraphics{../Figures/complexe.2}
\end{center}
\end{figure}





\`A tout nombre complexe $z=a+ib$ on associe le point $M$ de coordonnées
$(a,b)$ qu'on  appelle \textbf{image} de  complexe $z=a+ib$. On  le note
souvent $M(z)$. 

Inversement, à tout point $M$ du plan de coordonnées $(a,b)$, on associe
son \textbf{affixe} $z=a+ib$ qu'on note souvent $z_M$. 

Enfin,  à  tout   vecteur  $\ve{u}=a\ve{e_1}+b\ve{e_2}$  de  coordonnées
$(a,b)$  dans la  base $\pa{\ve{e_1},\ve{e_2}}$  est associé  une affixe
$z_{\ve{u}}=a+ib$ 



\subsection{Premiers calculs géométriques}

\begin{itemize}

\item[$\bullet$]  Soient   $\ve{u}$  et  $\ve{v}$  deux   vecteurs  de  coordonnées
  respectives         $(a,b)$         et        $(a',b')$,         alors
  $\ve{u}+\ve{v}=(a+a')\ve{e_1}+(b+b')\ve{e_2}$, donc 

  
    \begin{prop}[affixe d'une somme]
      \[z_{\ve{u}+\ve{v}}=z_{\ve{u}}+z_{\ve{v}}\]
      \end{prop}
    




\item[$\bullet$] De même, si $\la$ est un nombre réel 

  
    \begin{prop}[affixe du produit par un réel]
      \[z_{\la \ve{u}}=\la z_{\ve{u}}\]
      \end{prop}
    


\item[$\bullet$] Alors, si  $I$ est le \textbf{milieu} du segment $[A,B]$, on a


  
    \begin{prop}[affixe du milieu]
      \[z_I=\ud{(z_A+z_B)}\]
      \end{prop}
    




\item[$\bullet$] Pour tous points $A$ et $B$ 

  
    \begin{prop}[affixe d'un vecteur ]
      \[z_{\ve{AB}}=z_B-z_A\]
     \end{prop}
    

\end{itemize}


\subsection{Conjugué d'un complexe}

  \begin{defn}[Conjugué]
    On appelle conjugué du nombre complexe $z=a+ib$ le nombre $$\cacheg{\overline{z}=a-ib}$$
    \end{defn}
  



Géométriquement cela donne 



\begin{figure}
\begin{center}
  \includegraphics{../Figures/complexe.3}
\end{center}
\end{figure}



Je vous laisse prouver les propriétés immédiates suivantes~:


  \begin{prop}[]
    \begin{itemize}
\item[$\bullet$] $M(z)$ et $M'(\overline{z})$ sont symétriques par rapport à l'axe $(O,\ve{e_1})$
\item[$\bullet$] $\overline{z_1+z_2}=\cacheg{\overline{z_1}+\overline{z_2}}$
\item[$\bullet$] $\overline{z_1z_2}=\cacheg{\overline{z_1}\ \, \overline{z_2}}$
\item[$\bullet$] Soit $n\in \mathbb{Z}, z\in \mathbb{C}^*$, alors $\overline{z^n}=\overline{z}^n$
\item[$\bullet$] $\overline{\overline{z}}=\cacheg{z}$
\item[$\bullet$] $z\in\bbr\Longleftrightarrow z=\cacheg{\overline{z}}$
\item[$\bullet$] $z\in i\bbr\Longleftrightarrow z=\cacheg{-\overline{z}}$
\item[$\bullet$] $\RR e(z)=\cacheg{\ud{(z+\overline{z})}}$
\item[$\bullet$] $\IR m(z)=\cacheg{\ud{(z-\overline{z})}}$
\item[$\bullet$] Si $z=a+ib$, alors $z\overline{z}=\cacheg{a^2+b^2}$
\end{itemize}
    \end{prop}

\begin{demo}[$\overline{z_1z_2}=\cacheg{\overline{z_1}\ \, \overline{z_2}}$]
Soit $z_1=x_1+iy_1$ et $z_2=x_2+iy_2$ donc $z_1z_2=x_1x_2-y_1y_2+i(x_1y_2 + x_2y_1)$ alors $\overline{z_1z_2}=\overline{x_1x_2-y_1y_2+i(x_1y_2 + x_2y_1)}=x_1x_2-y_1y_2-i(x_1y_2 + x_2y_1)$.\\

D'autre part $\overline{z_1}\; \overline{z_2}=(x_1-iy_1)(x_2-iy_2)=x_1x_2-y_1y_2+i(-x_1y_2-x_2y_1)=x_1x_2-y_1y_2-i(x_1y_2 + x_2y_1)$.\\

Donc $$\overline{z_1z_2}=\overline{z_1} \; \overline{z_2}$$
\end{demo}  





\subsection{\`A quoi servent les conjugués ?}


\subsubsection{\`A montrer qu'un complexe est un réel}

En  effet, si  on arrive  à montrer  que $\overline{z}=z$,  alors  on en
conclut que $z$ est réel.\textbf{ Pourquoi~?}

\subsubsection{\`A rendre réel des dénominateurs pour obtenir des formes  algébriques} 

En
effet, $$z\cdot\overline{z}=(a+ib)(a-ib)=\cacheb{a^2-(ib)^2=a^2+b^2}$$ 


  \begin{exemple}[]
    Ainsi, pour obtenir la forme algébrique de l'inverse de $2+i$~: 

$\fr{1}{2+i}=\cacheb{\fr{1}{2+i}\times\fr{2-i}{2-i}=\fr{2+i}{4+1}=\fr{2}{5}+\fr{1}{5}i}$ 
    \end{exemple}
  

\subsection{Conjugué de l'inverse}


Sachant qu'un complexe non nul $z$ admet une forme algébrique $a+ib$, on
sait maintenant trouver la forme algébrique de son inverse~:
$$\dfrac{1}{z}=\dfrac{1}{a+ib}\times\dfrac{a-ib}{a-ib}=\dfrac{a-ib}{a^2-b^2}=\dfrac{a}{a^2-b^2}-i \dfrac{b}{a^2-b^2}$$
et donc $\overline{\left(\dfrac{1}{z}\right)}=\dfrac{z}{z \overline{z}}=\dfrac{1}{\overline{z}}$

\begin{demo}[$n\in \mathbb{Z}, \; \overline{z^n}=\overline{z}^n$ par "récurrence"]

Soit $z\neq 0$

Initialisation : $\overline{z^0}=\overline{1}=1$ et $ \overline{z}^0=1$


Hérédité : supposons que $\overline{z^n}=\overline{z}^n$ alors $\overline{z^{n+1}}=\overline{z^n z}=\overline{z^n} \overline{z}$ d'après la démonstration précédante.
Donc $\overline{z^{n+1}}=\overline{z^n}\overline{z}=\overline{z}^n \overline{z}=\overline{z}^{n+1}$



Conclusion: pour tout $n \in \mathbb{N}$, on a $  \overline{z^n}=\overline{z}^n$.

On prolonge le résultat sur $\mathbb{Z}$ de la manière suivante on pose $Z=1/z$
alors $\overline{Z^n}=\overline{Z}^n$ .

Puisque $\overline{\dfrac{1}{z}}=\dfrac{1}{\overline{z}}$ alors $\overline{\dfrac{1}{z^n}}=\dfrac{1}{\overline{z}^n}$ 

\end{demo}

\subsection{Module d'un nombre complexe}


  \begin{defn}[Module]
    Le module du complexe  $z$ est le réel positif noté $|z|$ tel que $$|z|=\sqrt{z\ \,\overline{z}}$$
    \end{defn}
  


    \begin{rem}[]
      \begin{itemize}
\item[$\bullet$] Cette définition en est bien une car $z\ \,\overline{z}=a^2+b^2$ d'après notre étude sur les conjugués. 
\item[$\bullet$] Si $a$ est un réel, $|a|=\sqrt{a\ \,\overline{a}}=\cacheg{\sqrt{aa}=\sqrt{a^2}}$ car $\overline{a}=a$. Donc le module de $a$ est bien la valeur absolue de $a$ et notre notation est cohérente.

La notion de module dans $\bbc$ généralise donc celle de valeur absolue dans $\bbr$.
\end{itemize}
      \end{rem}
    


\subsubsection{Interprétation géométrique}


\begin{figure}
\begin{center}
  \includegraphics{../Figures/complexe.4}
\end{center}
\end{figure}




Nous venons de voir que, si $z=a+ib$, alors 



  \begin{prop}[]
    \[|z|=\sqrt{a^2+b^2}\]
    \end{prop}
  


Or,  qu'est-ce que  $\sqrt{a^2+b^2}$ si  ce  n'est la  norme du  vecteur
$\ve{OM}$ ou encore la longueur $OM$. 



  \begin{prop}[]
    \[\ab{z_M}=\nov{OM}=OM\qquad \ab{z_{\ve{u}}}=\nov{u}\]
    \end{prop}
  



\subsubsection{Propriétés des modules}

Je vous laisse prouver les propriétés suivantes





  \begin{prop}[]
    \begin{itemize}
\item[$\bullet$] $\ab{\overline{z}}=\cacheg{\ab{z}}$
\item[$\bullet$] $\ab{z}=0\ssi \cacheg{ z=0}$
\item[$\bullet$] Soit $n\in\mathbb{Z}$ alors $\ab{z^n}=\ab{z}^n$
\item[$\bullet$] $\ab{z_1\cdot z_2}=\cacheg{\ab{z_1}\cdot\ab{z_2}}$
\item[$\bullet$] $\left|\frac{z_1}{z_2}\right|=\dfrac{\ab{z_1}}{\ab{z_2}}$
\item[$\bullet$] $\RR e(z)\ie \cacheg{\ab{z}}$
\item[$\bullet$] $\IR m(z)\ie \cacheg{\ab{z}}$
\end{itemize}
    \end{prop}
  
\begin{demo}[$n\in\mathbb{Z}$, $\ab{z^n}=\ab{z}^n$ par récurrence]
Initialisation : $\ab{z^0}=\ab{1}=1$. 

Hérédité : supposons que pour $n \in \mathbb{N}$ fixé on a $\ab{z^n}=\ab{z}^n$.
Ainsi $\ab{z^{n+1}}=\ab{z^n z}=\ab{z^n} \ab{z} = \ab{z}^n \ab{z} =\ab{z}^{n+1}$

Conclusion :  pour tout $n\in\mathbb{N}$, $\ab{z^n}=\ab{z}^n$

On utilise l'astuce $Z=\dfrac{1}{z}$ pour étendre le résultat sur $\mathbb{Z}$.
\end{demo}

La propriété suivante mérite une petite aide à la démonstration


  \begin{prop}[Inégalité triangulaire]
    $$\ab{z_1+z_2}\ie\ab{z_1}+\ab{z_2}$$
    \end{prop}
  



C'est  à dire, pour  aller de  Dijon à  Chalon-sur-Saône, il  est plus long de
passer par Paris que de suivre l'A31.  

Pour les curieux, voici comment cela se démontre.


Comme les deux membres de l'inégalité sont positifs, il suffit donc de comparer les carrés de chaque membre.

Or
$\ab{z_1+z_2}^2=\pa{z_1+z_2}\pa{\overline{z_1+z_2}}=(z_1+z_2)\pa{\overline{z_1}+\overline{z_2}}=\ab{z_1}^2+\pa{z_1\overline{z_2} 
+\overline{z_1}z_2}+\ab{z_2}^2$


D'autre part $\pa{\ab{z_1}+\ab{z_2}}^2=\ab{z_1}^2+2\ab{z_1z_2}+\ab{z_2}^2$

Il s'agit donc de comparer les <<~doubles produits~>>.

Or                                                     $z_1\overline{z_2}
+\overline{z_1}z_2=z_1\overline{z_2}+\overline{z_1\overline{z_2}}=2\RR
e\pa{z_1z_2}\ie 2\ab{z_1z_2}$ d'après une propriété ci-dessus. Donc 
$$\ab{z_1+z_2}^2=\ab{z_1}^2+\pa{z_1\overline{z_2}
+\overline{z_1}z_2}+\ab{z_2}^2\ie\ab{z_1}^2+2\ab{z_1z_2}+\ab{z_2}^2=\pa{\ab{z_1}+\ab{z_2}}^2$$




\section{R\'esolution d'\'equations du second degr\'e}

\subsection{Racine carrée d'un nombre complexe}

L'objet de cette section est de résoudre dans $\bbc$ l'équation $z^2=\al$

\subsubsection{Racine carrée d'un nombre réel}

On suppose ici que $\al$ est un réel.


\begin{itemize}
\item[$\bullet$] \boldmath $\al\se 0$\unboldmath : alors $z^2=\al\ssi 
  z^2-\al=$\cacher{$(z-\sqrt{\al})(z+\sqrt{\al})=0$}.  Les 
  solutions\footnote{LA solution si $\al=0$} 
  sont donc \cacher{$\pm \sqrt{\al}$}


    \begin{exemple}[]
      On connaît : $z^2=4\ssi $\cacheb{$z=-2$} ou \cacheb{$z=2$}
      \end{exemple}
    



\item[$\bullet$]    \boldmath   $\al<    0$\unboldmath    :   alors    $z^2=\al\ssi
  (z-i\sqrt{-\al})(z+i\sqrt{-\al})=0$.  Les  solutions  sont  donc  $\pm
  i\sqrt{-\al}$ 


  
    \begin{exemple}[]
      C'est la nouveauté : $z^2=-4\ssi$ \cacheb{$z=-2i$} ou \cacheb{$z=2i$}
      \end{exemple}
    

\end{itemize}


\subsubsection{Racine carrée d'un complexe non réel}

Les choses se  compliquent ! Nous allons traiter un  exemple pour ne pas
vous faire (trop) peur.  


  \begin{exemple}[]
    
Cherchons les racines carrées de $4+3i$, à savoir les nombres
$a+ib$ tels que $(a+ib)^2=a^2-b^2+2iab=4+3i$. 
    \end{exemple}
  

 Par unicité de la forme algébrique on obtient 

$\bsyc a^2-b^2&=&4\\a^2+b^2&=&5\\2ab&=&3\esyc$

Ainsi $a^2=$\cacher{$9/2$} et $b^2=$\cacher{$1/2$}, donc $a=$\cacher{$\pm3\sqrt{2}/2$} et
$b=$\cacher{$\pm\sqrt{2}/2$}, or $2ab=3$, donc $a$ et $b$ sont \cacher{de même signe}.

Les solutions sont donc \cacher{$\fr{\sqrt{2}}{2}(3+i)$} et
\cacher{$-\fr{\sqrt{2}}{2}(3+i)$}



\subsection{Résolution de ax\up{2}+bx+c=0 avec a, b   et  c  des réels} 

C'est        comme       en        première      :       \[ax^2+bx+c=0\ssi
a\Big[\big(x+\fr{b}{2a}\big)^2-\fr{b^2-4ac}{4a^2}\Big]=0\ssi
\bigl(x+\fr{b}{2a}\bigr)^2=\fr{b^2-4ac}{(2a)^2}\] 

Tout dépend donc du signe de $b^2-4ac$, puis on utilise les résultats de la section précédente.



  \begin{thm}[Résolution  de \boldmath$ax^2+bx+c=0\quad \text{avec}
    \ a,\ b \ \text{ et }\ c$ \unboldmath des réels] 
    L'équation $ax^2+bx+c=0$ admet toujours des solutions sur $\bbc$. 
Notons  $\Delta=b^2-4ac$  le   \textbf{discriminant}  de  l'équation  et
$\delta$ le \emph{complexe} vérifiant $\delta^2=\Delta$.
\begin{itemize}
\item[$\bullet$] Si $\Delta=0$, il existe une unique solution \cacheg{$x=-\fr{b}{2a}$}
\item[$\bullet$] Si $\Delta>0$, il existe deux solutions réelles  \cacheg{$x=\fr{-b\pm\sqrt{\Delta}}{2a}$}
\item[$\bullet$]  Si  $\Delta<0$, il  existe  deux  solutions complexes  conjuguées
  \cacheg{$x=\fr{-b\pm i\sqrt{-\Delta}}{2a}$} 
\end{itemize}
    \end{thm}
  



\section{Forme trigonométrique}

\subsection{Argument d'un complexe non nul}

\subsubsection{Forme trigonométrique}

Vous vous  souvenez de la correspondance  entre $\bbc$ et  le Plan. Nous
avions privilégié les coordonnées  cartésiennes d'un point. On aurait pu
utiliser tout  aussi bien ses  \textbf{coordonnées polaires}. Le  Plan a
cette fois besoin  d'être orienté (il le sera  implicitement à partir de
maintenant). 



\begin{figure}
\begin{center}
  \includegraphics{../Figures/complexe.5}
\end{center}
\end{figure}



Ainsi, $(r,\theta)$  étant le couple de coordonnées  polaires de l'image
$M$  de  $z$, on  a  $z=r\cos\theta+ir\sin\theta$  déterminé de  manière
unique, car c'est  en fait une forme algébrique  déguisée : on l'appelle
\textbf{forme trigonométrique} du complexe $z$. 



  \begin{defn}[forme trigonométrique]
    \[z=r\pa{\cos\theta+i\sin\theta}\]
    \end{defn}
  





\subsubsection{Congruence }

Vous  rencontrerez souvent  la  notation $x\equiv  y[2\pi]$  qui se  lit
<<~$x$ est congru à $y$  modulo 2$\pi$~>>. Elle veut simplement dire que
$x-y$ est  un multiple de  $2\pi$, c'est-à-dire qu'il existe  un entier
relatif $k$ tel que $x-y=k\times 2\pi$. 



  \begin{rem}[congruence modulo $2\pi$]
    \[x\equiv y[2\pi]\ssi \text{ il existe } k\in\bbz \text{ tel que } x=y+2k\pi\]
    \end{rem}
  


Par      exemple,       vous      savez      que      $\fr{\pi}{3}\equiv
\fr{4\pi}{3}[2\pi]$~:~dessinez un cercle 
trigonométrique pour vous en convaincre.

\subsubsection{Mesure d'un angle de vecteurs}

Nous n'avons  pas les moyens  de définir <<~proprement~>> les  angles de
vecteurs.  Nous n'en  avons  qu'une définition  intuitive.  Ce qui  nous
intéresse, c'est  que $\theta$ est UNE  mesure en radians  de l'angle de
vecteurs  $\pa{\ve{e_1},\ve{OM}}$.  UNE  mesure,  car elle  est  définie
modulo 2$\pi$.  Et bien cette mesure  sera UN argument  du complexe $z$,
qu'on notera $\arg z$. 
On retiendra~:


  \begin{prop}[argument]
    \[\arg z\equiv \theta[2\pi]\]
    \end{prop}
  



Par exemple, $\arg 32\equiv 0[2\pi]$, $\arg 32 i\equiv \fr{\pi}{2}[2\pi]$.

\subsubsection{Des formes trigonométriques de référence}

\begin{itemize}

\item[$\bullet$] $1=\cos 0+i\sin 0$  donc \cacher{$|1|=1$ et $\arg(1)\equiv 0[2\pi]$}
\item[$\bullet$]         $i=\cos\pa{\fr{\pi}{2}}+i\sin\pa{\fr{\pi}{2}}$        donc
  \cacher{$|i|=1$ et $\arg(i)\equiv\pa{\fr{\pi}{2}}[2\pi]$} 
\item[$\bullet$]                         $|1+i|=\sqrt{2}$                        et
  $1+i=\sqrt{2}\pa{\fr{\sqrt{2}}{2}+i\fr{\sqrt{2}}{2}}=$\cacher{$\sqrt{2}\pa{\cos\pa{\fr{\pi}{4}}+i\sin\pa{\fr{\pi}{4}}}$} 
donc \cacher{$\arg(1+i)\equiv\pa{\fr{\pi}{4}}[2\pi]$} 
\item[$\bullet$]    $|\sqrt{3}+i|=2$    et    $\sqrt{3}+i=2\pa{\fr{\sqrt{3}}{2}+\ud
    i}=$\cacher{$2\pa{\cos\pa{\fr{\pi}{6}}+i\sin\pa{\fr{\pi}{6}}}$} 
donc \cacher{$\arg(\sqrt{3}+i)\equiv\pa{\fr{\pi}{6}}[2\pi]$} 
\end{itemize}



\subsection{Correspondance forme algébrique - forme trigonométrique} 


Soit  $z\in\bbc$  de   forme  algébrique  $a+ib$  et   de  forme
trigonométrique $r\pa{\cos\theta+i\sin\theta}$ alors on a d'une part~: 


  \begin{prop}[forme algébrique connaissant la forme trigonométrique]
$a=$\cacheg{$r\cos\theta$}$\qquad b=$\cacheg{$r\sin\theta$}    
    \end{prop}
  



et d'autre part $r=|z|=\sqrt{a^2+b^2}$.



Si $z$ est \textit{non nul},  son module $r=\sqrt{a^2+b^2}$ sera non nul
également. Ainsi, nous pouvons écrire $z$ sous la forme 

\begin{center}
$\begin{array}{ccc}
z&=&\sqrt{a^2+b^2}\pa{\fr{a}{\sqrt{a^2+b^2}}+\I\fr{b}{\sqrt{a^2+b^2}}}\\
&=& r\pa{\fr{a}{\sqrt{a^2+b^2}}+\I\fr{b}{\sqrt{a^2+b^2}}}\\
   &=& r\pa{\cos(\theta)+\I\sin(\theta)}
\end{array}$
\end{center}

Nous en déduisons que~:



  \begin{prop}[forme trigonométrique en fonction de la forme algébrique]
$\cos(\theta)=$\cacheg{$\fr{a}{\sqrt{a^2+b^2}}$}$\qquad
\sin(\theta)=$\cacheg{$\fr{b}{\sqrt{a^2+b^2}}$} 
    \end{prop}
  




Ainsi, connaissant $a$ et $b$, on  peut obtenir le module et un argument
de  $a+\I   b$.  On   obtiendra  une  mesure   exacte  de   $\theta$  si
$\cos(\theta)$  et   $\sin(\theta)$  sont  des   valeurs  connues  comme
$\ofr{1}{2}$, $\ofr{\sqrt{3}}{2}$, 1, etc. 




  \begin{prop}[argument en fonction de la forme algébrique]
$\tan(\theta)=\fr{\sin(\theta)}{\cos(\theta)}=\fr{\fr{b}{\sqrt{a^2+b^2}}}{\fr{a}{\sqrt{a^2+b^2}}}=$\cacheg{$\fr{b}{a}$} 
    \end{prop}
  




ce qui déterminera une valeur de l'argument modulo $\pi$. 

Il  suffira ensuite  de  considérer  le signe  de  $\cos(\theta)$ ou  de
$\sin(\theta)$ pour savoir à qui on a affaire. 
















\subsection{Opérations sur les formes trigonométriques}

Soit $z=r\pa{\cos\theta+i\sin\theta}$ et $z'=r'\pa{\cos\theta'+i\sin\theta'}$,  alors 
$$zz'=rr'\cro{\pa{\cos\theta\cos\theta'-\sin\theta\sin\theta'}+i\pa{\sin\theta\cos\theta'+\cos\theta\sin\theta'}}$$ 
Vous qui connaissez parfaitement vos formules d'addition, vous en déduisez que 
$$zz'=z=rr'\pa{\cos\pa{\theta+\theta'}+i\sin\pa{\theta+\theta}'}$$

Ainsi, nous arrivons au résultat capital 


  \begin{prop}[argument d'un produit]
    $\arg(zz')=$\cacheg{$\arg(z)+\arg(z') \,[2\pi]$}
    \end{prop}
  




Cela va VOUS permettre de démontrer les propriétés suivantes avec un peu d'astuce et de patience



  \begin{prop}[Propriétés algébriques des arguments]
    \begin{itemize}
\item[$\bullet$] $\arg(zz')=$\cacheg{$\arg(z)+\arg(z') \,[2\pi]$}
\item[$\bullet$] $\arg(z^n)=$\cacheg{$n\arg(z)\,[2\pi]$}
\item[$\bullet$] $\arg\pa{\fr{1}{z}}=$\cacheg{$-\arg(z)\,[2\pi]$}
\item[$\bullet$] $\arg\pa{\fr{z}{z'}}=$\cacheg{$\arg(z)-\arg(z')\,[2\pi]$}
\item[$\bullet$] $\arg(\overline{z})=$\cacheg{$-\arg(z)\,[2\pi]$}
\item[$\bullet$] $\arg(-z)=$\cacheg{$\pi+\arg(z)\,[2\pi]$}
\end{itemize}
    \end{prop}
  


En particulier, la formule concernant $z^n$ nous permet d'écrire



  \begin{thm}[Formule de Moivre]
    \begin{center}
     $\pa{\cos \theta+i\sin\theta}^n= $\cacheg{$\cos(n\theta)+i\sin(n\theta)$}
    \end{center}
    \end{thm}
  



Nous nous rendons ainsi compte que~:


  \begin{rem}[]
    \begin{itemize}
\item[$\bullet$]  Les   formes  trigonométriques  sont  adaptés   aux  produits  de
  complexes~;
\item[$\bullet$] Les formes algébriques sont adaptées aux sommes de complexes.
  \end{itemize}
    \end{rem}
  






Vous  serez amenés  au hasard  d'exercices  à résoudre  des équations  à
coefficients  complexes, mais  on  n'attend de  vous aucun  savoir-faire
particulier : vous serez guidés pas à pas. 


\section{Les objets géométriques et les complexes}

%\section{De l'objet au complexe}

\subsection{Comment caract\'eriser un cercle ?}

Le cercle de centre $A$ et de rayon $R$ est l'ensemble des points situés
à la distance  $R$ de $A$. Il est facile de  traduire simplement cela en
langage complexe... 


  \begin{rem}[caractérisation d'un cercle]
    \[M(z)\in\CR(A,R)\ssi \ab{z-z_A}=R\]
    \end{rem}
  




\subsection{Comment caractériser un triangle isocèle ?}

C'est  encore une  histoire  de distance,  donc  de module  : $ABC$  est
isocèle de sommet principal $A$ si et seulement si $AB=AC$ donc 



  \begin{rem}[triangle isocèle]
    $ABC$     isocèle     de      sommet     principal     $A$     $\ssi
    \ab{z_B-z_A}=\ab{z_C-z_A}\ssi \ab{\fr{z_B-z_A}{z_C-z_A}}=1$ 
    \end{rem}
  



\subsection{Comment caractériser un triangle rectangle ?}

On peut encore parler distance grâce au théorème de Pythagore 
\[|z_C-z_B|^2=|z_B-z_A|^2+|z_C-z_A|^2\]

ou angle droit : mais c'est l'angle géométrique qui nous intéresse, donc
nous travaillerons modulo $\pi$ 
$$\pa{\ve{AB},\ve{AC}}=\pa{\ve{AB},\ve{e_1}}+\pa{\ve{e_1},\ve{AC}}=-\pa{\ve{e_1},\ve{AB}}+\pa{\ve{e_1},\ve{AC}}=-\arg\pa{z_{\ve{AB}}}+\arg\pa{z_{\ve{AC}}}=\arg\pa{\fr{z_{\ve{AC}}}{z_{\ve{AB}}}}=\arg\pa{\fr{z_C-z_A}{z_B-z_A}} $$ 



  \begin{rem}[triangle rectangle]
    $\pa{\ve{AB},\ve{AC}}\equiv\fr{\pi}{2}\,[\pi]\ssi
  \arg\pa{\fr{z_C-z_A}{z_B-z_A}}\equiv\fr{\pi}{2}\,[2\pi]$ 
    \end{rem}
  


\newpage

\subsection{Comment caractériser les différents quadrilatères ?} 


Petite révision de collège...




\begin{center}
  \includegraphics[angle=90,scale=0.7]{../Figures/complexe.7}
\end{center}



qu'il vous suffira d'adapter connaissant ce qui précède.
\newpage 
%\subsubsection{Comment caractériser une droite ?}





%\subsubsection{Comment caractériser certaines applications du plan ?}




\section{Du complexe à l'objet}




\subsection{Que représente $z-32+5i$ ?}

 Soit $A$ le point d'affixe $32-5i$ et $M$ le point d'affixe $z$, alors
 $z-32+5i=z_M-z_A=z_{\ve{AM}}$



\subsection{Comment interpréter $\vert z-32+5i \vert =3 $ ?}

D'après ce qui précède, on aboutit à $AM=3$ : il s'agit donc du cercle de centre $A$ et de rayon 3.

\subsection{Comment interpréter $ \vert 32+iz \vert =5 $ ?}

 $|32+iz|=|i(-32i+z)|=|i|\times|z-32i|=|z-32i|=BM$ avec $M$ le
 point d'affixe $z$ et $B$ le point d'affixe $32i$. On retombe donc sur un cercle.


\subsection{Comment interpréter $ \vert z-a \vert = \vert z-b \vert $ ?}

Soit $M$ d'affixe $z$, $A$ d'affixe $a$ et $B$ d'affixe $b$. Alors
l'égalité se traduit par $AM=BM$, donc $M$ est équidistant de $A$
et $B$, donc $M$ est sur la médiatrice de $[AB]$.

\subsection{Que     se     cache-t-il     derrière    le     quotient  $ \frac{z_C-z_A}{z_B-z_A} $ ?} 

Il             suffit            de             remarquer            que
$\fr{z_C-z_A}{z_B-z_A}=\fr{z_{\ve{AC}}}{z_{\ve{AB}}}$.     Donc     vous
utiliserez le fait que \begin{itemize} 
\item[$\bullet$] $\arg(\fr{z_C-z_A}{z_B-z_A})=(\ve{e_1},\ve{AC})-(\ve{e_1},\ve{AB})=(\ve{AB},\ve{AC})$
\item[$\bullet$] $\left|\fr{z_C-z_A}{z_B-z_A}\right|=\fr{AC}{AB}$
\end{itemize}



Dans  d'autres  cas,  vous  serez confrontés  à  l'interprétation  d'une
égalité du style $\fr{z_C-z_A}{z_B-z_A}=\la$ qui se traduit par 
$z_{\ve{AC}}=\la z_{\ve{AB}}$, donc \begin{itemize} 
\item[$\bullet$] si  $\la\in\bbr$, alors $\ve{AC}=\la\ve{AB}$  et donc $A$,  $B$ et
  $C$ sont alignés.  
\item[$\bullet$]  si  $\la\in  i\bbr$, $z_{\ve{AC}}=\pm|\la|iz_{\ve{AB}}$  et  donc
  $\arg\pa{z_{\ve{AC}}}\equiv\pm\fr{\pi}{2}+\arg\pa{z_{\ve{AC}}}\,[2\pi]$
  c'est à dire $(AC)\bot(AB)$ 
\item[$\bullet$] si $\la=\pm i$, alors le triangle $ABC$ est isocèle et rectangle en $A$ 
%\item[$\bullet$] si $\la=e^{\pm i\pi/3}$, alors le triangle ABC est équilatéral

\end{itemize}



\subsection{Comment interpréter qu'un angle de vecteurs est droit~? } 
%$\pa{\ve{MA},\ve{MB}} $}%\

On  déduit de cette relation que le triangle
$AMB$ est  rectangle en $M$, donc  que $M$ décrit le  cercle de diamètre
$[AB]$ privé des points $A$ et $B$. 


\subsection{Rotation dans le plan  $ \leftrightharpoons $ Multiplication}


Soit $z=3+2i$, alors $-1\times z=-3-2i$ et $i\times
z=3i+2i^2=-2+3i$. Notons $M$, $M'$ et $M''$ les points
d'\textbf{affixes} respectives $z$, $-z$ et $iz$



%\begin{figure}
%\begin{center}
 % \includegraphics{../Figures/complexe.8}
%\end{center}
%\end{figure}


Une multiplication par $i$ se traduit par un
quart de tour, une multiplication par $-1$ se traduit par un
demi-tour, deux multiplications successives par $i$, c'est à dire
une multiplication par $i^2=-1$ se traduit bien par deux quarts de
tour, $i.e.$ un demi tour.
 \begin{center}
\includegraphics{../Figures/complexe.8} 
\end{center} 
 
\section{Approche géométrique}

% \pro 
Si je vous dis transformation du plan, vous pensez sûrement à 
% \tsx
% 
des points qui se  transforment en d'autres points et qui gardent
plein de propriétés intéressantes : les longueurs, les angles, 
les  alignements,  le parallélisme  se  conservent.  Il  y a  juste  une
exception avec l'homothétie qui multiplie les longueurs. 

% \pro
C'est en  effet un résumé de vos  aventures géométriques des années
passées. Depuis, vous avez rencontré, lors de l'étude des nombres 
complexes  notamment,   d'autres  transformations  plus   étranges,  qui
transforment des droites en cercles ou en tout autre chose d'ailleurs.  

Nous allons nous occuper aujourd'hui d'une catégorie bien particulière de transformations, parmi bien
d'autres,  mais   avant  il  faudrait  s'entendre  sur   ce  qu'est  une
transformation du plan... 


\begin{defn}[transformation du plan]
  On dit que $f$ est une \textbf{transformation du plan} si 
\begin{itemize}
\item[$\bullet$] tout point du plan a une unique image par f
\item[$\bullet$] tout point du plan admet un unique antécédent par f
\end{itemize}
\end{defn}


Par exemple, une projection orthogonale n'est pas une transformation du plan selon notre définition.

Un exemple très important en géométrie, que vous connaissez depuis tout
petit (agrandissement-réduction), est l'homothétie~:

\begin{defn}[homothétie]
  Soit C un point du plan et $\lambda$ un réel non nul. \\ On appelle \textbf{homothétie de centre C et de rapport
$\lambda$} la transformation qui à tout point M du plan associe le point M$'$ tel que $$\ve{CM'}=\lambda\ve{CM}$$
\end{defn}


\begin{center}
  \includegraphics{../Figures/simi.1}
\end{center}


Un autre cas simple est la translation~:


\begin{defn}[translation]
   La \textbf{translation} de vecteur $\ve{v}$ est la transformation qui à tout point $M$ associe le point $M'$ tel que
$$\ve{MM'}=\ve{v}$$
\end{defn}


\begin{center}
  \includegraphics{../Figures/simi.2}
\end{center}


Nous étudierons enfin les rotations~:



\begin{defn}[rotation]
   La \textbf{rotation} de centre C et d'angle $\al$ est la transformation qui, à tout point M du plan, associe le point M$'$ tel
que \begin{itemize}
\item[$\bullet$] si M$\neq M'$, alors $CM=CM'$ et $(\ve{CM},\ve{CM'})=\al$
\item[$\bullet$] si $M=C$, alors $M'=C$
\end{itemize}
\end{defn}



\begin{center}
  \includegraphics{../Figures/simi.3}
\end{center}


\newpage
\section{Forme exponentielle d'un nombre complexe}

Avant d'aller plus loin, nous allons avoir besoin d'un petit outil technique. Notons $f$ la fonction définie sur $\bbr$ par $f(x)=\cos x+\mathrm{i}\sin x$. C'est une fonction un peu spéciale puisqu'elle est définie sur $\bbr$ mais est à valeurs dans $\bbc$.\\

 Que vaut $f(0)$~? $f(0)=\cos 0+\mathrm{i}\sin 0=1$ : jusqu'ici, tout va bien.
Je vous demande à présent un petit effort d'imagination : on peut supposer que $f$ est dérivable sur $\bbr$ en considérant $i$ comme
un coefficient quelconque et en extrapolant les formules de dérivations des fonctions à valeurs réelles. Qu'est-ce que ça peut donner ? $f'(x)=-\sin x+\mathrm{i}\cos x$, le lien avec $f(x)$ apparaît puisque $f'(x)=\mathrm{i}(\mathrm{i}\sin x+\cos x)=\mathrm{i} f(x)$\\ 

Récapitulons : $f$ vérifie $f'=\mathrm{i} f$ avec $f(0)=1$. Ça ne vous rappelle rien~? 

La fonction exponentielle ! Quand on avait $f'=kf$ et $f(0)=1$ alors on en concluait que $f(x)=\mathrm{e}^{kx}$. Donc vous ne  serez pas choqué si nous écrivons, par convention  d'écriture à notre niveau, que $f(x)=\exp(\mathrm{i} x)=\mathrm{e}^{\mathrm{i} x}$. 


\begin{thm}[notation exponentielle]
  On convient d'écrire, pour tout réel $x$ $$\exp^{\mathrm{i} x}=\cos x+\mathrm{i}\sin x$$
\end{thm}

\begin{thm}[Formules d'euler]
Soit $z=e{i\theta}$ alors on a en faisant $\dfrac{z+\overline{z}}{2}$ et $\dfrac{z+\overline{z}}{2i}$, \\
$$\cos \theta = \dfrac{e^{i \theta} +e^{-i \theta}} {2}$$
$$\sin \theta = \dfrac{e^{i \theta} -e^{-i \theta}} {2i}$$
\end{thm}


Notez  au passage  que les  formules d'additions  sont  alors plus  faciles à  retrouver. En  effet,
$\mathrm{e}^{\mathrm{i}(a+b)}=\mathrm{e}^{\mathrm{i} a}\mathrm{e}^{\mathrm{i} b}$, donc

\begin{eqnarray*}
  \cos(a+b)+\mathrm{i}\sin(a+b)&=&(\cos a+\mathrm{i}\sin a)(\cos b+\mathrm{i}\sin b)\\
&=&(\cos a\cos b-\sin a\sin b)+\mathrm{i}(\cos a\sin b+\cos b\sin a)
\end{eqnarray*}


 Alors, par unicité de
la forme algébrique d'un complexe, on obtient~:
$$\cos(a+b)=\cos a\cos b-\sin a\sin b\quad \text{et} \quad \sin(a+b)=\cos a\sin b+\cos b\sin a$$

Voici une illustration du côté extrêmement pratique de l'outil complexe.



\section{Écriture complexe des transformations usuelles}

\subsection{Translations}

Considérons la translation de vecteur $\ve{v}$ d'affixe $b$ et les habituels points $M$ et $M'$ d'affixes $z$ et $z'$. Par
définition, on  a $\ve{MM'}=\ve{v}$, à l'aide des nombres complexe on trouve
$z_{\ve{MM'}}=z'-z=b$, i.e. $z'=z+b$. Ici on retrouve une expression du type $az+b$, donc c'est bien une
similitude directe  de rapport $|a|=|1|=1$,  donc une isométrie.

\begin{prop}[translation : version complexe]
 La transformation complexe définie par  $$z\mapsto z+b$$ est la translation de vecteur d'affixe $b$
\end{prop}


\subsection{Rotations}

\begin{prop}[rotation : version complexe]
  La transformation complexe définie par $$z\mapsto ze^{i\al}+b$$ est une rotation d'angle $\al$ avec $\al\not\equiv0[2\pi]$
\end{prop}


\begin{center}
\includegraphics{../Figures/simi.4}
\end{center}

\subsection{Homothéties}

Avec les  notations habituelles, on obtient $\ve{CM'}=k\ve{CM}$, d'où $z'-c=k(z-c)$, i.e.
$z'=kz+c(1-k)$. Ainsi une homothétie a une représentation complexe de la forme $z\mapsto kz+b$. l faut faire attention : on ne doit pas avoir $z=z+b$ qui nous fait retomber sur une translation, donc $k$ doit être différent de 1.


\begin{prop}[homothéties : version complexe]
  La transformation complexe définie par $$z\mapsto kz+b\quad \text{avec}\quad k\in\bbr\setminus\{0,1\}$$ est une homothétie
de rapport $k$
\end{prop}
 

\chapter{Les équations polynomiales}
Dans tout ce cours, $K$ d\'{e}signe $\mathbb{R}$ ou $\mathbb{C}$%
\section{Définitions des  Fonctions polyn\^{o}mes.}

\begin{defn}

une application $f$ d'une partie $I$ de $K$ dans $K$ est dite \textit{%
polynomiale} (ou appel\'{e}e une \textit{fonction polyn\^{o}me}) si 
\begin{equation*}
\exists n\in \mathbb{N}\;\exists \left( a_{0},a_{1},...,a_{n}\right) \in
K^{n+1}\;/\;\forall x\in I\;\;f\left( x\right) =\underset{k=0}{\overset{n}{%
\sum }}a_{k}x^{k}=a_{0}+a_{1}x+...+a_{n}x^{n}
\end{equation*}

L'ensemble des fonctions polynomiales de $I$ dans $K$ est not\'{e} $\mathcal{%
P}\left( I,K\right) $. Cette ensemble à le bon goût d'être stable par multiplication et addition de polynôme.

\end{defn}


\begin{defn}

\qquad - le \textit{degr\'{e}} d'un polyn\^{o}me non nul est l'indice
maximum d'un coefficient non nul ; par convention, le degr\'{e} du polyn\^{o}%
me nul est $-\infty .$%
\begin{equation*}
\text{pour }P=\left( a_{k}\right) _{k\geqslant 0},\deg P=\underset{a_{k}\neq
0}{\max }k
\end{equation*}

\bigskip \qquad - la \textit{valuation} d'un polyn\^{o}me non nul est
l'indice minimum d'un coefficient non nul ; par convention, la valuation du
polyn\^{o}me nul est $+\infty .$%

\end{defn}

%\begin{equation*}
%\text{pour }P=\left( a_{k}\right) _{k\geqslant 0}, \limfunc{val}P=\underset{%
%a_{k}\neq 0}{\min }k
%\end{equation*}

\begin{defn}
\qquad - les polyn\^{o}mes de degr\'{e} 0 et le polyn\^{o}me nul sont dits 
\textit{constants}.

%\qquad - $P$ est appel\'{e} un \textit{mon\^{o}me} si $\deg P=\limfunc{val}P$
(un seul coefficient non nul).

\qquad - si $n=\deg P,$ le coefficient de rang $n$ est appel\'{e} le
coefficient \textit{dominant} ou "\textit{de t\^{e}te}" du polyn\^{o}me.

\qquad - un polyn\^{o}me dont le coefficient dominant est \'{e}gal \`{a} $1$
est dit \textit{normalis\'{e}}, ou \textit{unitaire}.
\end{defn}

\begin{prop}

\begin{equation*}
\forall P,Q\in K[X]\;\;\forall x\in K\;\left\{ 
\begin{array}{c}
\left( P+Q\right) \left( x\right) =P\left( x\right) +Q\left( x\right) \\ 
\left( PQ\right) \left( x\right) =P\left( x\right) Q\left( x\right) \\ 
\left( P\circ Q\right) \left( x\right) [\text{ou }\left( P\left( Q\right)
\right) \left( x\right) ]=P\left( Q\left( x\right) \right)%
\end{array}
\right.
\end{equation*}

\end{prop}
\section{Division euclidienne des polyn\^{o}mes.}

\bigskip

\begin{thm}[Division euclidienne des polynomes]

 \'{e}tant donn\'{e} deux polyn\^{o}mes $A,B\neq 0,$ il existe un unique
couple $\left( Q,R\right) $ de polyn\^{o}mes v\'{e}rifiant 
\begin{equation*}
A=BQ+R,\text{ avec }\deg R<\deg B
\end{equation*}

$Q$ et $R$ sont appel\'{e}s respectivement le \textit{quotient} et le 
\textit{reste} de la division euclidienne de $A$ par $B.$

\end{thm}

\section{Polyn\^{o}mes premiers entre eux, PGCD, PPCM.}

\begin{defn}
Il existe un unique polynôme $D$ nul ou unitaire tel que, pour tout $R$ de $K[X]$, on ait :
($R\vert P$ et $R\vert Q$) si et seulement si $R\vert D$.
Ce polynôme D est appelé le pgcd de $P$ et $Q$, on le note aussi $PGCD(P,Q)$ ou $P\wedge Q$.\\

$P$ et $Q$ sont premiers entre eux ssi leur $PGCD$ est $1$.
\end{defn}

\begin{thm}[Théorème de B\'{e}zout ]

$A$ et $B$ sont premiers entre eux si et seulement s'il
existe deux polyn\^{o}mes $U$ et $V$ tels que $AU+BV=1.$

\end{thm}

\begin{cor}[Théorème de Gauss]
Si $A$ divise $BC$ et $A$ et $B$ sont
premiers entre eux, alors $A$ divise $C$.
\end{cor}


\begin{rem}
On peut poser $A\wedge 0=0\wedge A=A.$
\end{rem}

\begin{defn}
Soient $P$ et $Q$ deux éléments de $K[X]$. Il existe un unique polynôme $M$ nul ou unitaire tel que, pour tout $R$ de $K[X]$, on ait :
($P \vert R$) et $Q\vert R$) si et seulement si $M \vert R$.
Ce polynôme $M$ est appelé le ppcm de $P$ et $Q$.
\end{defn}


\begin{defn}
$PPCM\left( A,B\right) .PGCD\left( A,B\right) =AB$ si $A$ et 
$B$ sont ont des coefficients dominants égaux à 1.
\vspace{5mm}
\end{defn}

\section{Racines des polynôme}


\subsection{D\'{e}finition et premiers exemples}

\begin{defn} soit $P\in K[X]$ et $x_{0}\in K$ ; on dit que $x_{0}$ est une racine
(ou un z\'{e}ro) de $P$ si $P\left( x_{0}\right) =0$.
\vspace{5mm}
\end{defn}

\begin{thm}[Racine et divisibilité] 
$x_{0}$ est une racine de $P$ si et seulement si le polyn\^{o}me $X-x_{0}$ divise $P$
dans $K[X]$.
\end{thm}

\begin{demo}
Faisons la division euclidienne de $P$ par $(X-x_0)$. Ainsi ils existent $R,Q \in K[X]$ tels que $P=(X-x_0)Q+R$ avec $deg(R)<1$ donc $deg(R)=0$ autrement dit $R$ est un polynôme constant.\\

On sait que $P(x_0)=0$ donc $(X-x_0)Q+R)(x_0)=0$, autrement dit $R(x_0)=0$.$R$ étant constant, on déduit que est le polynôme nul, d'où $R=0$.\\

On conclut, $P=(X-x_0)Q$ ce qui par définition équivaut à dire que le polyn\^{o}me $X-x_{0}$ divise $P$.
\end{demo}
\subsection{Formule fondamentale de l'algèbre}
\begin{thm}
$$X^n-a^n=(X-a)(X^{n-1}+X^{n-2}a+\cdots+Xa^{n-2}+a^{n-1})$$\\
Autrement dit,
$$X^n-a^n=(X-a)\sum_{k=1}^n-1 X^{n-k}a^k$$
\end{thm}

\begin{demo}
Soit $P=X^{n-1}+X^{n-2}a+\cdots+Xa^{n-2}+a^{n-1}$, alors $XP=X^n+X^{n-1}a+\cdots+Xa^{n_1}$ et $aP=X^{n-1}a+\cdots+Xa^{n_1}+a^n$. \\
Ainsi $XP-aP=X^n-a^n$, et $XP-aP=(X-a)P=(X-a)((X^{n-1}+X^{n-2}a+\cdots+Xa^{n-2}+a^{n-1})$
\end{demo}




\section{Nombre de racines d'un polyn\^{o}me.}


\begin{prop}
un polyn\^{o}me non nul a toujours un nombre fini de racines
distinctes, inf\'{e}rieur ou \'{e}gal \`{a} son degr\'{e}.
\end{prop}

\begin{demo}
Supposons par l'absurde qu'il existe un polynôme $P$ de degré $n$ qui possèdent $k$ racines avec $k>n$ notées $\al_1, \al_2, \cdots, \al_k$ alors $P=(X-\al_1)(X-\al_2)\cdots(X-\al_k)Q$ avec $deg(Q) \se0 $.\\
 
Or $deg(P)=deg((X-\al_1)(X-\al_2)\cdots(X-\al_k)Q)=deg((X-\al_1))+deg((X-\al_2))+\cdots+deg((X-\al_k))+deg(Q) =k+deg(Q)$.\\

Donc d'un coté $deg(P)=n$ et d'autre part $deg(P)=k+deg(Q)(\se k)$, donc $n\se k$ or $k>n$ par hypothèse, d'où l'absurdité.
\end{demo}

\section{Racines $n$-ième et polygones réguliers}

\subsection{Résolution générale de l'équation $z^n-a=0$}
Soit $a\in \mathbb{C}$, on pose $a=r_ae^{i\theta_a}$ son écriture exponentielle avec $r_a \se 0$ et $\theta_a$ un argument de $a$ compris entre $[0;2\pi[$. De même supposons que $z=re^{i\theta}$ soit solution avec $r \se 0$ et $\theta \in [0;2\pi[$
alors on a nécéssairement 
$$r^ne^{in\theta}=r_ae^{in\theta}$$ 
En appliquant le module à cette équation on déduit facilement que $r=\sqrt[n]{r_a}$.\\

En appliquant l'argument à cette équation on déduit aussi nécéssairement que $$n\theta \equiv \theta_a \pmod {2\pi}$$
autrement dit, 
$$\theta \equiv \dfrac{\theta_a}{n} \pmod {\dfrac{2\pi}{n}}$$

il existe donc $n$ arguments possibles sur un segment de longueur $2\pi$.\\

Ainsi pour tout $k\in\{0,1,\cdots,n-1\}$, $z_k=\sqrt[n]{r_a}e^{  \dfrac{\theta_a}{n} +\dfrac{2k\pi}{n} } $ est une racine du polynôme $z^n-a=0$. Or un polynôme de degré $n$ à au plus $n$ racines, donc le polynôme $z^n-a$ à exactement $n$ racine qui sont les $(z^k)^{n-1}_{k=0}$. \\

Puisque $z^n-a$ est divisible par tous les $z-z_k$ on a $z^n-a=\prod^{n-1}_{k=0}(z-z_k)Q$ avec $deg(Q)=0$. \\

\textit{Effectivement $z^n-a=(z-z_0)Q_0$ pour un certain $Q_0$ alors on sait que $deg(z^n-a)=deg((z-z_0)Q_0)=deg((z-z_0))+deg(Q_0)=1+deg(Q_0)$ donc $n=1+deg(Q_0)$ d'où $deg(Q_0)=n-1$.}
\textit{De plus $z^n-a(z_1)=(z-z_0)Q_0(z_1)$ ce qui équivaut à dire que $0=(z_1-z_0)Q_0(z_1)$ et comme $(z_1-z_0)\neq 0$ alors nécéssairement $Q_0(z_1)=0$ ainsi $Q_0=(z-z_1)Q_1$ avec $deg(Q_1)=deg(Q_0)-deg(z-z_1)=n-1-1=n-2$.}
\textit{En répétant ce procédé $n$ fois on obtient le résultat $z^n-a=\prod^{n-1}_{k=0}(z-z_k)Q$ avec $deg(Q)=0$}

Reste donc à déterminer le polynôme constant $Q$, le coefficient dominant de $z^n-a$ est 1. \\

Le coefficient dominant de $\prod^{n-1}_{k=0}(z-z_k)Q$ est le produit des coefficients dominants de chaque termes et donc égale au coefficient dominant de $Q$ qui est la constante recherchée.\\

Par l'égalité $z^n-a=\prod^{n-1}_{k=0}(z-z_k)Q$ on déduit l'égalité des coefficients dominants et donc $Q=1$ nécéssairement. \\

On a donc la formule suivante. 

$$z^n-a= \prod^{n-1}_{k=0}(z-z_k) =\prod^{n-1}_{k=0}(z-\sqrt[n]{r_a}e^{(  
\frac{\theta_a}{n} +\frac{2k\pi}{n} )} )$$

\subsection{Liens avec polygones réguliers}
On va représenter nos racines dans le plan complexe.\\

Une des premières choses à remarquer est que le module des racines est constante, ainsi toutes les racines sont situés sur le cercle $\calc (O,\sqrt[n]{r_a})$, ensuite on remarque que l'angle formé $\widehat{Z_kOZ_{k+1}}$ est toujours de $\dfrac{2\pi}{n}$ \emph{de même pour $\widehat{Z_{n+1}OZ_0}$}.\\

Ainsi, on $n$ affixes dans le plan,  tous sont sur le cercle $\calc (O,\sqrt[n]{r_a})$ espacées d'un angle régulier  $\dfrac{2\pi}{n}$. Reste à déterminer un seul point pour fixé le polygone, prenons une affixe $Z_k$ quelconque et le tour est joué.


\subsection{Cas particuliers: Racine  $n$-ième de l'unité}
Soit $n\se 1$ un entier. On appelle racine $n$-ième de l'unité dans  $\bbc$, tout élément $z$ de $\bbc$ vérifiant $z^n=1$. Une racine $n$-ième de l'unité est donc une racine du polynôme $X^n-1$ de $\bbc [X]$.\par
Si on cherche $z$ sous la forme exponentielle $z=\rho\exp^{\imath \theta}$, on voit que $\rho^n\exp^{\imath n\theta}=1$, ce qui donne $\rho=1$ et $\theta=2k\pi/n$ avec $0\ie k<n$.\\

 Les racines  $n$-ièmes de l'unité dans $\bbc$ sont les $n$ nombres complexes

 $\exp^{2\imath k\dfrac{\pi}{n}}=\cos(2k\dfrac{\pi}{n})+\imath\sin(2k\dfrac{\pi}{n})$,  avec $0\ie k<n$, on note $\mathbb{U}_n$ l'ensemble des racines $n$-ièmes de l'unité. \\
 

 
Pour $n=2$, les racines carrées de l'unité sont 1 et $-1=\exp^{\imath\pi}$.\\

Pour $n=3$, les racines cubiques de l'unité sont 1, $\jmath$ et $\jmath^2$ avec $\jmath=\exp^{2\imath\pi/3}=-\dfrac{1}{2}+\imath\dfrac{\sqrt3}{2}$ et $\jmath^2=\exp^{4\imath\pi/3}=-\dfrac{1}{2}-\imath\dfrac{\sqrt3}{2}$.\\
\begin{center}
\includegraphics{../Figures/polyn.5}
\end{center}
Pour $n=4$, les racines quatrièmes de l'unité sont $\pm1$, $\imath=\exp^{\imath\pi/2}$ et $-\imath=\exp^{3\imath\pi/2}$.\\
\begin{center}
\includegraphics{../Figures/polyn.4}
\end{center}
Pour finir, voici deux autres illustrations pour n=5 et n=6, etc.

\begin{minipage}{7cm}
\includegraphics[scale=0.7]{../Figures/polyn.6}
\end{minipage}
\hfill
\begin{minipage}{7cm}
\includegraphics[scale=0.7]{../Figures/polyn.1}
\end{minipage}


Comme $X^n-1=(X-1)(\sum_{1\ie k\ie n}X^{n-k})$, toute racine $n$-ième de l'unité distincte de 1 est racine du polynôme $\sum_{1\ie k\ie n}X^{n-k}=X^{n-1}+X^{n-2}+\ldots+1$. Autrement dit, si $\zeta$ est une racine $n$-ième de l'unité, elle vérifie les deux relations $\zeta^n=1$ et, si $\zeta\not=1$, $\sum_{1\ie k\ie n}\zeta^{n-k}=0$ qui sont très utiles dans certains calculs.\\

Pour calculer avec des racines $n$-ièmes de l'unité, on peut utiliser leur forme exponentielle $\exp^{2\imath k\pi/n}$ ou leur forme trigonométrique $\cos(2k\pi/n)+\imath\sin(2k\pi/n)$.


\begin{rem}
Une subtilité et non pas des moindres est à remarquer ici, l'ensemble des racines $n$-ième de l'unité $\mathbb{U}_n$ à le bon gout s'il est muni de la multiplication complexe, d'être stable par $\times$, de contenir un élément neutre $1$,  et qui de plus est stable par inverse.\\

Un tel ensemble est qualifié comme groupe, une notion très importante en mathématiques.

\end{rem}
\section{Dérivation des Polynômes, formules de Taylor}

\section{D\'{e}finition de l'application dérivation}

\begin{defn}
le polyn\^{o}me d\'{e}riv\'{e} du polyn\^{o}me $P=\underset{k\geqslant
0}{\sum }a_{k}X^{k}$ est le polyn\^{o}me, not\'{e} $P^{\prime }=\underset{%
k\geqslant 1}{\sum }ka_{k}X^{k-1}=\underset{k\geqslant 0}{\sum }\left(
k+1\right) a_{k+1}X^{k}.$\\

Les polyn\^{o}mes d\'{e}riv\'{e}s successifs se notent de la m\^{e}%
me fa\c{c}on que pour les fonctions.

On notera $D$ l'application : $\left\{ 
\begin{array}{c}
K[X]\rightarrow K[X] \\ 
P\mapsto P^{\prime }%
\end{array}%
\right. $
\end{defn}

\begin{prop}

P1 : $P\in K\Leftrightarrow P^{\prime }=0$

\bigskip

P2 : si $\deg \left( P\right) \geqslant 1,\;\deg $ $P^{\prime }=\deg P-1,$
et mieux, si $n\geqslant 1,$ $P\in K_{n}[X]\Leftrightarrow P^{\prime }\in
K_{n-1}[X]$

\bigskip

P3 : $\left( P+Q\right) ^{\prime }=P^{\prime }+Q^{\prime }$ ; $\left(
\lambda P\right) ^{\prime }=\lambda P^{\prime }$ ; $\left( PQ\right)
^{\prime }=P^{\prime }Q+PQ^{\prime }$

\bigskip

P4 : $\left( P\circ Q\right) ^{\prime }=\left( P^{\prime }\circ Q\right)
Q^{\prime }$

\end{prop}

\section{Formule de Taylor}

\section{Formule de Taylor pour les polyn\^{o}mes}

\bigskip
\begin{prop}
Si deg$P=n,$%
\begin{equation*}
P=P\left( X\right) =\underset{k=0}{\overset{n}{\sum }}\dfrac{P^{\left(
k\right) }\left( x_{0}\right) }{k!}\left( X-x_{0}\right) ^{k}=P\left(
x_{0}\right) +P^{\prime }\left( x_{0}\right) \left( X-x_{0}\right)
+P^{\prime \prime }\left( x_{0}\right) \dfrac{\left( X-x_{0}\right) ^{2}}{2}%
+..+P^{\left( n\right) }\left( x_{0}\right) \dfrac{\left( X-x_{0}\right) ^{n}%
}{n!}
\end{equation*}

ou, ce qui revient au m\^{e}me 
\begin{equation*}
P\left( x_{0}+X\right) =\underset{k=0}{\overset{n}{\sum }}\dfrac{P^{\left(
k\right) }\left( x_{0}\right) }{k!}X^{k}=P\left( x_{0}\right) +P^{\prime
}\left( x_{0}\right) X+P^{\prime \prime }\left( x_{0}\right) \dfrac{X^{2}}{2}%
+..+P^{\left( n\right) }\left( x_{0}\right) \dfrac{X^{n}}{n!}
\end{equation*}

\end{prop}

\begin{cor}

un polyn\^{o}me de degr\'{e} $n$ est enti\`{e}rement d\'{e}termin\'{e} par
la connaissance de $P\left( x_{0}\right) ,P^{\prime }\left( x_{0}\right)
,P^{\prime \prime }\left( x_{0}\right) ,...,P^{\left( n\right) }\left(
x_{0}\right) $
\end{cor}

\chapter{Arithmétique}

\section{Divisibilité dans $\mathbb{Z}$ et congruences}

 \begin{defn}
Soit $a$ et $b$ deux entiers relatifs.
On dit que $a$ \underline{divise} $b$ s'il existe un entier $k$ tel que $b=a\times k$.
Lorsque $a$ divise $b$, on note $a\mid b$ 

On dit aussi que: 
\begin{tabular}[t]{lll}$a$ est un \textbf{diviseur} de $b$\\$b$ est \textbf{divisible} par $a$\\$b$ est un    \textbf{multiple} de $a$
\end{tabular}

\end{defn}

 \begin{rem}

$0$ est un multiple de tout entier: pour tout $n\in \mathbb{Z}, 0 \times n=0$


 \end{rem}

\begin{prop}
Soit $a$ et $b$ deux entiers

a) Si $a \mid b$ et $b \neq 0 $ alors $ \,\vert a \vert \leqslant \vert b \vert$

b) Tout entier $b$ \textbf{non nul }a un nombre fini de diviseurs.

\end{prop}

 \begin{demo}
a) Comme $a \mid b$, on peut écrire $b=k \times a$ avec $k\in \mathbb{Z}$ donc:

$\vert b \vert=\vert k \times a \vert = \vert k\vert \times \vert a\vert$
mais $b\neq 0$ donc $k \neq 0$ et $\vert k\vert \geq1$ et par conséquent:$\vert a\vert\leq\vert b\vert$

b) Soit $b$ un entier non nul. Si $a$ est un diviseur de $b$ on a vu que: $\vert a\vert\leq\vert b\vert$
donc $-\vert b \vert\leq a\leq \vert b \vert$ et $a$ peut prendre au maximum $2 \times \vert b \vert$ valeurs.

En revanche, 0 a une infinité de diviseurs car tous les entiers divisent 0.

\end{demo}


 \begin{prop} 
Soit $a, b$ et $c$ trois entiers.

a) Si $a\mid b$ et $b\mid c$ alors $a\mid c$

b) $c\mid a$ si et seulement si $c\mid -a$

{\small (L'ensemble des diviseurs de $a$ est égal à l'ensemble des diviseurs de $-a$)}



c)Si $c\mid a$ alors $c\mid ab$

d) Si $c\mid a$ et $c\mid b$ alors: \begin{tabular}[t]{lll}1) $c\mid a+b$\\2) $c\mid a-b$\\3) $c\mid au+bv$ avec $u$ et $v$ entiers
\end{tabular}

e) Soit $a$ et $b$ non nuls. Si $b\mid a$ et $a \mid b$ alors $a=b$ ou $a= -b$

\end{prop}

 \begin{demo}

b) Si $c\mid a$, il existe un entier $k$ tel que: $a=k \times c$ et donc $-a=\left( -k\right) \times c$ soit $c\mid-a$

Si $c\mid -a$, il existe un entier $k'$ tel que: $-a=k' \times c$ et donc $a=\left( -k'\right) \times c$ soit $c\mid a$

\end{demo}






\begin{thm}
Soit $a$ un entier et $b$ un naturel non nul.

Il existe un unique entier $q$ et un unique entier $r$ tels que:
\begin{center}
$a=b \times q+r$ avec $0\leq r < b$
\end{center}

\end{thm}

\begin{demo}
On considère les multiples de $b: ...-2b, -b, 0, b, 2b, ...$

L'entier $a$ est:

* soit multiple de $b$ et dans ce cas il existe un unique entier $q$ tel que: $a=b \times q$

* soit compris strictement entre deux multiples de $b$ et donc il existe un unique entier $q$ tel que:

\hspace{3cm}$bq < a < b\left( q+1\right) $

Si on note $r=a-bq$ on a bien, dans les deux cas, $a=b \times q+r$ avec $0\leq r < b$

\end{demo}


\begin{rem}

$a$ est le dividende,
$b$ est le diviseur,
$q$ est le quotient et
$r$ est le reste
\end{rem}

\begin{exemple}
Déterminer le quotient et le reste de la division euclidienne de $a$ par $b$ avec:

a) $a=325, b=7$

b) $a=-113, b=7$


Réponses: \\

a) $325=7 \times 46+3$

b) $113=7 \times 16 +1$ d'où $-113=-7 \times 16 -1$ puis $-113=-7 \times 17+6$, $q=-17, r=6$


\end{exemple}






\begin{defn}
Soit $a$ et $b$ deux entiers et $n$ un entier naturel non nul.

On dit que $a$ est \textbf{congru à} $b$ \textbf{modulo} $n$ si $a$ et $b$ ont le même reste dans la division euclidienne par $n$.

\end{defn}

\begin{rem}

$a$ est congru à $b$ modulo $n$ se note au choix:
$a \equiv b \left( n\right) $ ou $a \equiv b \left[ n\right] $ ou $a \equiv b$ mod $ n$
 
 \end{rem}
 
\begin{exemple}
86 et 23 ont pour reste 2 dans la division euclidienne par 7 donc $86\equiv 23 \left( 7\right) $

\end{exemple}

\begin{thm}
Soit $a$, $b$, $c$ et $r$ entiers et $n$ un entier naturel non nul.

1) $a \equiv b \left( n\right) $ si et seulement si $a-b$ est un multiple de $n$

2) Si $a \equiv r \left( n\right) $ et si $0 \leq r < n$ alors $r$ est le reste de la division de $a$ par $n$.

3) Si $a \equiv b \left( n\right) $ et  $b \equiv c \left( n\right) $ alors $a \equiv c \left( n\right) $

\end{thm}



\begin{demo}[ $a \equiv b \left( n\right) $ si et seulement si $a-b$ est un multiple de $n$]

  $\bigstar$ Si $a \equiv b \left( n\right) $ alors la division  par $n$ donne: $a=nq+r$ et $b=nq'+r$ donc $a-b=n(q-q')$: $n$ divise $a-b$
  
 $\bigstar$ Si $n$ divise $a-b$ il existe $k \in\mathbb{Z}$ tel que $a-b=nk$ soit $a=b+nk$
  
  La division de $a$ par $n$ se traduit par: $a=nq+r$ avec $0\leq r <n$. On a alors $b+nk=nq+r$ donc $b=n(q-k)+r$ et $0\leq r < n$ qui se traduit par: $a$ et $b$ ont le même reste dans la division par $n$

\end{demo}


\begin{prop}[Congruences et opérations]
Soit $a$, $b$, $c$ et $d$ entiers et $n$ un entier naturel non nul.

1) Si $a \equiv b \left( n\right) $ alors $ac \equiv bc \left( n\right) $

2) Si $a \equiv b \left( n\right) $ et $c \equiv d \left( n\right) $ alors:

\begin{tabular}[t]{lll}1) $a+c \equiv b+d \left( n\right) $\\2) $a-c \equiv b-d \left( n\right) $\\3) $ac \equiv bd \left( n\right) $
\end{tabular}

3) Si $a \equiv b \left( n\right) $ alors, pour tout naturel $p$, on a: $a^{p} \equiv b^{p} \left( n\right) $

\end{prop}

\section{PGCD}




\begin{defn}

Soit $a$ et $b$ deux entiers. On note $\Delta\left( a;b\right) $ l'ensemble des diviseurs communs à $a$ et $b$.
Si $b\mid a$ alors $\Delta\left( a;b\right) $ est égal à l'ensemble des diviseurs de $b$, qu'on note $\Delta\left( b\right) $

\end{defn}

\begin{prop}
Si $a, b, q$ et $r$ sont des entiers tels que: $a=bq+r$ alors $\Delta\left( a;b\right)=\Delta\left( b;r\right)  $
 \end{prop}
 
\begin{demo}

Si $d\in \Delta\left( a;b\right)$ alors $d\mid r$ car $r=a-bq$ donc $d\in \Delta\left( b;r\right)$

Si $d\in \Delta\left( b;r\right)$ alors $d\mid a$ car $a=bq+r$ donc $d\in \Delta\left( a;b\right)$


Remarque: cas particulier\\

Si $a\in \mathbb{Z}, b\in\mathbb{N}^{*}$ et $a=bq+r, 0 \leq r < b $ est la division euclidienne de $a$ par $b$ alors: $\Delta\left( a;b\right)=\Delta\left( b;r\right)  $


\end{demo}

\begin{defn}

Si $a$ et $b$ sont deux entiers non tous les deux nuls alors $ \Delta\left( a;b\right)$ est un ensemble fini et on appelle PGCD de $a$ et $b$ le \textbf{plus grand diviseur commun} de $a$ et $b$.
\end{defn}

\begin{rem}
pgcd$(a;b)$ ou $a\wedge b$
\end{rem}

\begin{prop}
Soit $a$, $b$ et $r$ entiers non nuls.

1) pgcd$(a;b)\geq 1$

2) pgcd$(a;b)$ = pgcd$(b;a)$

3) pgcd$(a;a)=\vert a\vert$

4) pgcd$(a;b)$ =pgcd$(-a;b)$ =pgcd$(a;-b)$

5) Si $b\mid a$, pgcd$(a;b)=\vert b \vert$

6) Si $a=bq+r$  alors,  pgcd$(a;b)$ =pgcd$(b;r)$ 


\end{prop}

\begin{exemple}[Principe de l'algorithme d'Euclide pour calculer le pgcd]

Soit $a$ et $b$ deux entiers naturels non nuls tels que $b$ ne divise pas $a$. Alors pgcd$(a;b)$ est le \textbf{dernier reste non nul} de la suite des divisions de l'algorithme d'Euclide.

Si $r_{1},r_{2},...,r_{n}$ est la suite de restes non nuls, on a:

$ \Delta\left( a;b\right)= \Delta\left( b;r_{1}\right)=... \Delta\left( r_{n-1};r_{n}\right)= \Delta\left( r_{n}\right)$


\end{exemple}

\begin{exemple}

Calculer pgcd$(8820;3150)$

\underline{Solution}
$8820=3150\times 2+2520$

$3150=2520\times 1+630$

$2520=630\times 4+0$ 
donc pgcd$(8820;3150)=630$

\end{exemple}

\begin{prop}

Soit $a$ et $b$ deux entiers non nuls et $d= pgcd(a;b)$

1) L'ensemble des diviseurs communs à $a$ et $b$ est l'ensemble des diviseurs de $d$

2) Il existe des entiers $u$ et $v$ tels que: $d=au+bv$

3) Pour  $k \in\mathbb{N}^{*}, pgcd(ka;kb)=k\times pgcd(a;b)$

\end{prop}

\begin{demo}

1)$\Delta(a;b)=\Delta(d)$ d'après l'algorithme d'Euclide.

2) Soit $E=\left\lbrace au+bv \text{ avec } u,v\in \mathbb{Z}\right\rbrace $

Si $u=1,v=0$, on voit que $a\in E$
Si $u=-1,v=0$, on voit que $-a\in E$
Si $u=0,v=1$, on voit que $b\in E$

$E$ contient au moins un entier strictement positif. Soit $\delta$ le plus petit d'entre eux;
il existe $u_{0}$ et $v_{0}$ entiers tels que: $\delta=au_{0}+bv_{0}$

Soit $au+bv$ un élément quelconque de $E$ La division euclidienne de cet élément par $\delta$ donne
  $au+bv=\delta q+r, 0 \leq r < \delta $ d'où: $r=au+bv-q \left(au_{0}+bv_{0}\right)=a \left( u-qu_{0}\right) +b\left( v-qv_{0}\right) $
 
 
 
 Par suite, $r\in E$ et $0 \leq r < \delta$, mais d'après la définition de $\delta$: $r=0$ donc \textbf{tout élément de E est multiple de } $\delta$.
 
 $\delta\mid a$ et $\delta\mid b$ donc $\delta\mid d=pgcd(a;b)$.
 
 Réciproquement, $d\mid a$ et $d\mid b$ donc aussi $d\mid au_{0}+bv_{0}$ c'est-à-dire, $d\mid\delta$. Conclusion: $d=\delta$
 
 
\end{demo}




\section{Entiers premiers entre eux}


\begin{defn}

Soit $a$ et $b$ deux entiers non nuls. On dit que $a$ et $b$ sont \textbf{premiers entre eux} si pgcd$(a,b)=1$

\end{defn} 
 
\begin{thm}

Soit $a,b$ et $d$ des entiers non nuls. 

pgcd$(a,b)=d$ si et seulement si, il existe deux entiers non nuls $a'$ et $b'$ tels que: $a=da',b=db'$ et $a'\wedge b'=1$
 
\end{thm}

\begin{demo}

  Si $d=$ pgcd$(a,b), d\mid a, d\mid b$ donc il existe $a', b' \in \mathbb{Z}^{*}$ tels que: $a=da', b=db'$
  
  Soit $r=$ pgcd$(a',b')$ alors pgcd$(da',db')=dr=d$ donc $dr=d$ avec $d\neq 0$ donc $r=1$ et $a'\wedge b'=1$
  
  Réciproquement, si $a=da', b=db'$ et $a'\wedge b'=1$ alors pgcd$(da',db')=d \times$pgcd$(a',b')=d$
 
\end{demo}


\begin{thm}[Théorème de Bézout]

Deux entiers non nuls $a$ et $b$ sont premiers entre eux si et seulement si, il existe des entiers $u$ et $v$ tels que: $au+bv=1$
\end{thm}

\begin{demo}

Si $a$ et $b$ premiers entre eux, pgcd$(a,b)=1$ et donc il existe deux entiers $u$ et $v$ tels que $au+bv=1$

Réciproquement, si $au+bv=1$ et $d\mid a$ et $d\mid b$ alors $d\mid au+bv$ donc $d$ divise 1. Les seuls diviseurs communs à $a$ et $b$ sont $-1$ et $1$ et $a \wedge b=1$

\end{demo}

\begin{cor}
Soit a,b deux entiers tel que $pgcd(a,b)=d$ alors il existe deux entiers $u$ et $v$ tels que,

 $$ua+vb=d$$
\end{cor}

\begin{demo}
Soit a,b deux entiers tel que $pgcd(a,b)=d$ alors $pgcd(\frac{a}{d},\frac{b}{d})=1$ par définition donc il existe deux entiers $u$ et $v$. Par le théorème de Bézout on a:
 $$u\dfrac{a}{d}+v\dfrac{b}{d}=1$$
 autrement dit 
 $$ua+vb=d$$

\end{demo}

\begin{thm}

Soit $a$, $b_{1}$ et $b_{2}$ des entiers non nuls.

Si $a\wedge b_{1}=1$ et $a\wedge b_{2}=1$ alors $a\wedge (b_{1}b_{2})=1$  
\end{thm}

\begin{demo}

Si $a\wedge b_{1}=1$ il existe $u,v$ tels que $au+b_{1}v=1$

Si $a\wedge b_{2}=1$ il existe $u',v'$ tels que $au'+b_{2}v'=1$

donc: $(au+b_{1}v)(au'+b_{2}v')=1=a(auu'+b_{2}uv'+b_{1}u'v)+(b_{1}b_{2})(vv')$ donc $a\wedge (b_{1}b_{2})=1$  


\end{demo}

\begin{thm}[Théorème de Gauss]

Soit $a,b,c$ entiers avec $a, b$ non nuls.

Si $a$ divise $bc$ et $a$ premier avec $b$ alors $a$ divise $c$.

\end{thm}

\begin{demo}

Comme $a \wedge b=1$ il existe $u,v$ tels que $au+bv=1$ donc $cau+cbv=c$

$a$ divise $acu$ et $bcv$ donc $a$ divise $c$.

\end{demo}

\begin{thm}

Un entier $n$ divisible par $a$ et $b$ tels que $a\wedge b=1$ est divisible par $ab$.
\end{thm}

\begin{demo}

Si $a\mid n$ il existe un entier $k$ tel que $n=ak$.

Si $b\mid n$ alors $b\mid ak$ et comme $a\wedge b=1$, d'après Gauss, $b\mid k$. Il existe donc $k'$ tel que $k=bk'$ et on a alors $n=ak=abk'$ donc $n$ est divisible par $ab$

\end{demo}


\section{Nombres premiers}


\begin{defn}

Un entier naturel est dit \textbf{premier} s'il admet exactement deux diviseurs positifs (1 et lui-même)

\end{defn}

\begin{rem}
1 n'est pas premier.
\end{rem}

\begin{thm}

 Tout $n\in \mathbb{N}$ avec $n\geq 2$ admet au moins un diviseur premier.

 Si $n$ n'est pas premier et $n\geq 2$ alors il admet un diviseur premier compris entre 2 et $\sqrt{n}$
 \end{thm}


\begin{demo}

Si $n$ est premier, il admet bien un diviseur premier: lui-même.

Si $n$ n'est pas premier alors il admet un plus petit diviseur positif $p\neq 1$.
$p$ est premier sinon $p$ aurait lui-même un diviseur positif différent de 1 qui serait un diviseur de $n$, mais plus petit que $p$.

De plus, $n$ peut s'écrire $n=p \times r$ avec $p\leq r$ donc $p^{2}\leq p \times r$ soit $p^{2} \leq n$ et $p\leq\sqrt{n}$

\end{demo}

\begin{rem}
On utilise ce théorème de la manière suivante:

Si un naturel $n\geq 2$ n'admet pas de diviseur premier compris entre 2 et $\sqrt{n}$ alors $n$ est premier.

\end{rem}

\begin{exemple}

Pour savoir si 631 est premier, il suffit de tester tous les nombres premiers inférieurs à $\sqrt{631}\approx25,12$

\end{exemple}

\begin{thm}
Il existe une infinité de nombres premiers.
\end{thm}

\begin{demo}

On suppose qu'il existe un nombre fini de nombres premiers $p_{1}, p_{2}, .... ,p_{n}$


Soit $N=p_{1}\times p_{2}\times .... \times p_{n}+1$

$N$ n'est pas premier car pour tout $i$ de 1 à $n$, $N> p_{i}$.

 $N$ admet donc un diviseur premier dans la liste $p_{1}, p_{2}, .... ,p_{n}$. 

Soit $p_{k}$ ce diviseur premier de $N$.

$p_{k}$ divise $N$ et $p_{k}$ divise $p_{1}\times p_{2}\times .... \times p_{n}$ donc

$p_{k}$ divise $N-p_{1}\times p_{2}\times .... \times p_{n}$, soit $p_{k}$ divise 1: la seule possibilité est que $p_{k}=1$ qui est impossible car $p_{k}$ est premier.

\end{demo}

\begin{thm}

Tout entier naturel strictement supérieur à 1 admet une décomposition, unique à l'ordre des facteurs près, en produit de nombres premiers.

\end{thm}

\begin{exemple}

$72=2^{3}\times 3^{2}$
\end{exemple}


\begin{thm}[Petit théorème de Fermat]
Si $p$ est un nombre premier et si $a$ est un entier quelconque, alors $a^p-a$ est un multiple de $p$.
\end{thm}

\begin{lemme}
$p\vert C^k_p$ pour $k\in \{1,2,3,\cdots, p-1 \}$\\
\end{lemme}

\begin{demo}

$$C^k_p=\dfrac{p!}{(p-k)!k!}$$ 
$$C^k_p(p-k)!k!=p! $$
$$C^k_p(p-k)!k!=p\times (p-1)!$$

\end{demo}

\begin{lemme}
Soit $p$ un nombre premier $(k+1)^p \equiv 1+k^p \pmod p$
\end{lemme}

\begin{demo}
Par la formule du binôme on a,
$$(k+1)^p=\sum_{i=0}^p C^i_pk^i$$
de plus pour $i \in \{1,2,\cdots,p-1\}$ on a $p\vert C^i_p$ donc $ C^i_p\equiv 0 \pmod p$. 
Donc 
$$(k+1)^p \equiv 1+k^p$$
\end{demo}

\begin{demo}[Petit théorème de Fermat, par récurrence]
Initialisation : Pour $a=1$ on a $a^p=1^p=1$ donc l'assertion est vrai. \\

Hérédité : Supposons que maintenant que l'assertion est vrai pour $a=k$ alors montrons que l'assertion est vraie au rang $k+1$.\\

$(k+1)^p \equiv 1+k^p \pmod p$ or $k^p \equiv k \pmod p$ par hypothèse de récurrence. Donc $(k+1)^p \equiv 1+k \pmod p$.\\

Conclusion: pour tout entier $k$ et $p$ un nombre premier $k^p=k$.

\end{demo}


\chapter{Les matrices}

 \section{Définitions}

\begin{defn}
Une \textbf{matrice} $A$ de dimension $n\times p$ ou de format $(n;p)$ est un tableau de nombres comportant $n$ lignes et $p$ colonnes.
   
   On note $a_{ij}$ l'élément se trouvant à l'intersection de la ligne $i$ et de la colonne $j$.
   
   Lorsque $n=p$ on dit que la matrice est une matrice \textbf{carrée} \textbf{d'ordre} $n$.
   
\end{defn} 

\begin{exemple}
   
    $A=\begin{pmatrix}
    2 & -5 & 4 \\ 
   6 & 3 & -8
   \end{pmatrix} $ est une matrice de format $(2;3)$
   
   Une matrice carrée d'ordre 2 est une matrice de la forme: $\begin{pmatrix}
    a_{11} & a_{12}  \\ 
   a_{21} & a_{22} 
   \end{pmatrix}$
   
\end{exemple}
   
\subsection{Matrices particulières}
   
    Si $n=1$, $A$ est une \textbf{matrice ligne}.
   
     Si $p=1$, $A$ est une \textbf{matrice colonne}.
   
    Si tous les coefficients sont nuls, $A$ est une \textbf{matrice nulle}.
   
\subsection{Égalité de matrices}
   \begin{defn}
   
    Deux matrices sont égales lorsqu'elles ont le même format et les mêmes coefficients aux mêmes emplacements.
   
   \end{defn}
   
\begin{defn}
   
   La matrice \textbf{transposée} d'une matrice $A$ de format $(n;p)$ est la matrice de format $(p;n)$, notée $A^{T}$ obtenue en échangeant les lignes et les colonnes de $A$.
   
\end{defn}   

\begin{exemple}

Si $A=\begin{pmatrix}
    2 & -5 & 4 \\ 
   6 & 3 & -8
   \end{pmatrix} $  alors $A^{T}=\begin{pmatrix}
    2 & 6 \\ 
   -5 & 3 \\
   4 &-8
   \end{pmatrix} $ 
   
\end{exemple}
   
 \subsection{Addition de deux matrices de même format}
   
    On appelle \textbf{somme } de deux matrices de même format la matrice obtenue en additionnant les coefficients de même emplacement.
   
\subsection{Multiplication d'une matrice par un nombre réel}
   
    Pour un réel $k$ et une matrice $A$, on note $kA$ la matrice $M$ dont l'élément $m_{ij}$ est égal à $ka_{ij}$
   
\begin{exemple}
   
   $2\begin{pmatrix}
    2 & -5 & 4 \\ 
   6 & 3 & -8
   \end{pmatrix} = \begin{pmatrix}
    2\times 2 & 2\times (-5) &2\times 4 \\ 
   2\times 6 & 2\times 3 & 2 \times (-8)
   \end{pmatrix}=\begin{pmatrix}
    4 & -10 & 8 \\ 
   12 & 6 & -16
   \end{pmatrix} $ 
   
  \end{exemple}
   
\begin{prop}
   
    $A,B$ et $C$ sont des matrices de même format, $O$ est la matrice nulle de même format, $k$ et $k'$ sont deux nombres réels.


   
   \begin{enumerate}
   \item $A+B=B+A$
   
   \item $A+(B+C)=(A+B)+C$
   
   \item $A+O=O+A=A$
   
   \item $0A=O$ et $1A=A$
   
   \item $(k+k')A=kA+k'A$ et $k(A+B)=kA+kB$ 
   \end{enumerate}
\end{prop}   

\subsection{Multiplication de deux matrices}
   
  \begin{defn} 
   $\blacktriangleright$ Le produit de la matrice ligne $L=\begin{pmatrix}
   a_{1}&a_{2}&...&a_{p}
   \end{pmatrix}$ par la matrice colonne $C= \begin{pmatrix}
   b_{1}\\
   b_{2}\\
   .\\
   .\\
   .\\
   b_{p}
   
   \end{pmatrix}$
      est le nombre $LC=a_{1}b_{1}+a_{2}b_{2}+...+a_{p}b_{p}=\displaystyle\sum_{k=1}^{p}a_{k}b_{k}$
      
   $\blacktriangleright$ Le \textbf{produit} de la matrice $A=\left( a_{ij}\right) $ de format $(n;p)$ par une matrice $B=\left( b_{ij}\right) $ de format $(p;r)$ est la matrice, notée $AB$, de format $(n;r)$ dont le coefficient $(c_{ij})$ est le produit de la matrice ligne $i$ de $A$ par la matrice colonne $j$ de $B$:
   
 \begin{center}
  \boxed{ $$c_{ij}=\displaystyle\sum_{k=1}^{p}a_{ik}b_{kj}$$}
 \end{center}
 \end{defn}
 
\begin{exemple}
 
 $A$ de format $(2;3)$ $B$ de format $(3;2)$, $C=AB$ de format $(2;2)$
  \begin{center}
 \begin{tabular}{cc}
   \begin{pspicture}(2,2)
  \psline{*->}(0.4,-0.6)(2.4,0.6)
  \psline{*->}(1,-0.6)(2.4,0.1)
  \psline{*->}(1.6,-0.6)(2.4,-0.4)
  \end{pspicture}  &$ \begin{pmatrix}
   1 & -1 \\ 
   3 & 2 \\
   -3 & 4 
   \end{pmatrix} $ \\ 
 $\begin{pmatrix}
    1 & 3 & 5 \\ 
   2 & 4 & 6
   \end{pmatrix} $ & $\begin{pmatrix}
    -5 & 25  \\ 
   -4 & 30 
   \end{pmatrix} $ \\ 
 
 \end{tabular} 
 \end{center}
 
 \end{exemple}

\begin{rem}

 
  Si les matrices $AB$ et $BA$ sont définies, $\boxed{ \text{en général } AB\neq BA}$.
  \end{rem}
  
  
\begin{prop}

   $A,B$ et $C$ sont des matrices dont les formats permettent les calculs indiqués, $k$ est un réel
  
  \begin{enumerate}
  \item $A(BC)=(AB)C$
  
  \item $A(B+C)=AB+AC$
  
  \item $(A+B)C=AC+BC$
  
  \item $(kA)B=A(kB)=k(AB)$
\end{enumerate}   

\end{prop}

\subsection{Matrices unités}
\begin{defn}
Soit $n$ un entier naturel non nul. On appelle \textbf{matrice unité d'ordre $n$} la matrice $I$, carrée d'ordre $n$ dont tous les coefficients sont nuls sauf ceux de la diagonale principale (éléments $a_{ii}$) qui sont égaux à 1.
\end{defn}

\begin{exemple}

La matrice unité d'ordre 3 est $I=\begin{pmatrix}
    1 & 0 & 0 \\ 
   0 & 1 & 0 \\
   0 & 0 & 1
   \end{pmatrix} $
  
\end{exemple}

   \section{Inverse d'une matrice carrée}
  
   \begin{defn}
   $A$ est une matrice carrée d'ordre $n$. On dit qu'une matrice $B$, carrée d'ordre $n$, est \textbf{l'inverse} de $A$ si elle vérifie $AB=I$ et $BA=I$.
	\end{defn}   
   
\begin{prop}
   
   Si la matrice carrée $A$ d'ordre $n$ admet une inverse, celle-ci est unique. On la note $A^{-1}$.
   
\end{prop}
   
 \section{Puissance d'une matrice carrée}
   \begin{defn}
   
   $A$ est une matrice carrée et $n \in \mathbb{N}^{*}$. La puissance n-ième de la matrice $A$, notée $A^{n}$, est la matrice définie par: $A^{n}=\underbrace{AA....A}_{\text{n fois}}$
   
    Par convention, $A^{0}=I$
   
  \end{defn}
  
   \begin{thm}
   
   Pour $n \in \mathbb{N}^{*}$ et pour tous nombres réels $a$ et $b$,
    
    \begin{center}
    $\begin{pmatrix}
    a & 0  \\ 
   0 & b 
   \end{pmatrix}^{n}=\begin{pmatrix}
    a^{n} & 0  \\ 
   0 & b^{n} 
   \end{pmatrix}$
    \end{center}
   \end{thm}
   
    \begin{demo}
    par récurrence: exercice
	\end{demo}
	
	    
\begin{thm}
    
    Soit $A=\begin{pmatrix}
    a & b  \\ 
   c & d 
   \end{pmatrix}$ une matrice carrée d'ordre 2.
   
   1) Si $ad-bc\neq 0$, $A$ admet une inverse $\boxed{A^{-1}=\dfrac{1}{ad-bc}\begin{pmatrix}
    d &-b  \\ 
   -c & a 
   \end{pmatrix}}$
   
   2) Si $ad-bc=0$, $A$ n'a pas d'inverse.
   
  \end{thm}
   
\begin{demo}
   
   1) Soit $B=\dfrac{1}{ad-bc}\begin{pmatrix}
    d &-b  \\ 
   -c & a 
   \end{pmatrix}$. On vérifie que $AB=BA=I=\begin{pmatrix}
    1 & 0  \\ 
   0 & 1 
   \end{pmatrix}$
   
   2)  On suppose que $A$ admet une inverse $A'$
   
   Soit $B=\begin{pmatrix}
    -c & a  \\ 
   -c & a 
   \end{pmatrix}$ et $C=\begin{pmatrix}
    d &-b  \\ 
   d & -b 
   \end{pmatrix}$
   
   $BA=\begin{pmatrix}
    -c & a  \\ 
   -c & a 
   \end{pmatrix}$
    $ \begin{pmatrix}
    a & b  \\ 
   c & d 
   \end{pmatrix}=\begin{pmatrix}
    0 & 0  \\ 
   0 & 0
   \end{pmatrix}=O$
   De même, $CA=O$
   
   $B=B(AA')=(BA)A'=O$ donc $a=c=0$
   
   $C=C(AA')=(CA)A'=O$ donc $d=b=0$, $A$ est la matrice $O$ 
   
   ce qui implique que: $AA'=I=O$ ce qui est impossible: $A$ n'a pas d'inverse.\\
  \end{demo} 
   
   \begin{rem}
    Le nombre $ad-bc$ s'appelle le \textbf{déterminant }de la matrice $A$\\
   \end{rem}
   
 \section{Application aux systèmes linéaires}
   
   Tout système linéaire de $n$ équations à $n$ inconnues peut s'écrire sous la forme matricielle $AX=B$ où $A$ est la matrice carrée d'ordre $n$ des coefficients du système, $X$ est la matrice colonne des inconnues et $B$ est la matrice colonne formée par les seconds membres des équations.
   
   Si la matrice carrée $A$ est inversible alors le système a une unique solution égale à $A^{-1}B$.
   
\section{Matrices et suites }
\subsection{Suites de matrices}

Une suite de matrices colonnes de taille $k$ ( $k \in \mathbb{N}, k\geq 2$) est une fonction de $\mathbb{N}$ dans l'ensemble des matrices colonnes de taille $k$.

\begin{defn}

On dit que la suite de matrices colonnes $(X_{n})$ de taille $k$ est \textbf{convergente} si les $k$ suites formées par les termes correspondant à la même ligne sont convergentes.

La limite de la suite est alors la matrice colonne formée des $k$ limites obtenues.
\end{defn}

\subsection{Suites de la forme $U_{n+1}=AU_{n}+B$}

\begin{prop}

Soit une suite de matrices colonnes $(U_{n})$ de taille $k$ telle que, pour tout $n \in \mathbb{N}$,

\begin{center}
$\boxed{U_{n+1}=AU_{n}, \text{ où } A \text{ est une matrice carrée d'ordre }k}$
\end{center}

Alors, pour tout $n$ de $\mathbb{N}$ , $\boxed{U_{n}=A^{n}U_{0}}$\\

\end{prop}

\begin{demo}
\vspace{3cm}
\end{demo}

\begin{prop}


Soit une suite de matrices colonnes $(U_{n})$ de taille $k$ vérifiant pour tout $n \in \mathbb{N}$ , $U_{n+1}=AU_{n}+B$, où $A$ est une matrice carrée non nulle d'ordre $k$ et $B$ une matrice colonne de taille $k$.

S'il existe une matrice $C$ telle que $C=AC+B$ alors le terme général de cette suite peut s'écrire:

\begin{center}
$\boxed{U_{n}=A^{n}(U_{0}-C)+C}$
\end{center}

\end{prop}

\begin{rem}
Si $I-A$ est inversible alors $C=(I-A)^{-1}B$
\end{rem}
 
\begin{demo}
 
 Comme  $U_{n+1}=AU_{n}+B$ et $C=AC+B$, par différence on a: $U_{n+1}-C=A(U_{n}-C)$
 
 La suite $(V_{n})$ telle que $V_{n}=U_{n}-C$ vérifie $V_{n+1}=AV_{n}$ donc $V_{n}=A^{n}V_{0}$
 
 D'où: $U_{n}-C=A^{n}(U_{0}-C)$ soit $U_{n}=A^{n}(U_{0}-C)+C$
 
\end{demo} 

\subsection{Convergence des suites vérifiant $U_{n+1}=AU_{n}+B$}
 
 Soit une suite de matrices colonnes vérifiant $U_{n+1}=AU_{n}+B$
 
 On suppose qu'il existe une matrice $C$ telle que $C=AC+B$
 
 1) Si $U_{0}=C$, la suite converge vers $C$
 
 2) Si $U_{0}\neq C$ et si la suite $(A^{n})$ converge vers une matrice $A'$ alors la suite $(U_{n})$ converge vers
 
 \begin{center}
 $A'(U_{0}-C)+C$
 \end{center}
 
 
\begin{demo}
\vspace{3cm}
\end{demo}
 

\section{Marche Aléatoire}

 \subsection{Marche aléatoire entre deux états}
    \begin{defn}
         
          On considère un système qui n'a que deux états possibles A et B et qui évolue par étapes successives.
           
           On note $p$ la probabilité qu'il passe de A à B.
           
           On note $q$ la probabilité qu'il passe de B à A.
           
           Le \textbf{graphe probabiliste} ci-contre donne l'évolution du système d'une étape à la suivante.
           
           \vspace{0.7cm}
          \begin{center}
           \begin{psmatrix}[mnode=circle]
              A &  & B 
           \end{psmatrix}
           \psset{arrows=->,shortput=nab}
           \nccircle[angleA=0]{1,1}{0.5cm}^{1-p}
           \ncarc[arcangle=15]{1,1}{1,3}^{p}
           \ncarc[arcangle=15]{1,3}{1,1}^{q}
           \nccircle[angleA=0]{1,3}{0.5cm}^{1-q}
          \end{center}
          
          \vspace{0.3cm}
          
          On définit la \textbf{matrice de transition} par: $T=\begin{pmatrix}
    1-p & p  \\ 
   q & 1-q
   \end{pmatrix}$
   
   \end{defn}
   
\begin{rem}
   
   Tous les coefficients appartiennent à $[0;1]$
   
   Pour chaque ligne, la somme des coefficients est 1.
   
   \end{rem}
   
  \begin{defn}
   
   Pour $n \in \mathbb{N}$ , on note:
   
   $A_{n}$ l'événement: "à l'étape $n$ le système est dans l'état A".
   
   $B_{n}$ l'événement: "à l'étape $n$ le système est dans l'état B".
   
   $a_{n}=p(A_{n}), b_{n}=p(B_{n})$. On a: $\boxed{a_{n}+b_{n}=1}$
\end{defn}   
   
\begin{defn}
   
   La matrice \textbf{ligne} $P_{n}=(a_{n}\hspace{0.3cm} b_{n})$ est appelée la \textbf{répartition de probabilité à l'étape $n$}.
   
  \end{defn}

\begin{prop}
   
   Pour $n \in \mathbb{N},\boxed{P_{n+1}=P_{n}T}$
   
   \end{prop}
   
   \begin{demo}
   
   
   \begin{tabular}{cc}
   \psset{nodesep=0mm,levelsep=20mm,treesep=10mm}
\pstree[treemode=R]{\Tdot}
{
\pstree
{\Tdot~[tnpos=a]{$A_{n}$}\taput{\small $a_{n}$}}
{
\Tdot~[tnpos=r]{$A_{n+1}$}\taput{\small $1-p$}
\Tdot~[tnpos=r]{$B_{n+1}$}\tbput{\small $p$}
}
\pstree
{\Tdot~[tnpos=a]{$B_{n}$}\tbput{\small $b_{n}$}}
{
\Tdot~[tnpos=r]{$A_{n+1}$}\taput{\small $q$}
\Tdot~[tnpos=r]{$B_{n+1}$}\tbput{\small $1-q$}
}
}
\begin{tabular}{rl}


 $p(A_{n+1})=$&$ p(A_{n}\cap A_{n+1})+p(B_{n}\cap A_{n+1})$\\
 =&$(1-p)a_{n}+qb_{n}$ \\
 -----------------&-----------------------------------------\\
 $p(B_{n+1})=$&$p(A_{n}\cap B_{n+1})+p(B_{n}\cap B_{n+1})$\\
 =&$pa_{n}+(1-q)b_{n}$
 \end{tabular}\\ 
 
   \end{tabular} 
   
  
\begin{center}
 \begin{tabular}{cc}
    &$ \begin{pmatrix}
    \hspace{1.1cm}1-p & p \textcolor{white}{.} \\ 
  \hspace{1.1cm} q & 1-q \textcolor{white}{.} 
   \end{pmatrix} $ \\ 
 $P_{n}T=\begin{pmatrix}
     a_{n} & b_{n}
   \end{pmatrix} $ & $\begin{pmatrix}
    (1-p)a_{n}+b_{n}q & pa_{n}+(1-q)b_{n}  
   
   \end{pmatrix} $ \\ 
 
 \end{tabular} 
 \end{center}

  On a bien:  $P_{n+1}=P_{n}T$
  
  \end{demo}
  
\begin{prop}

   Pour tout $n \in \mathbb{N}:P_{n}=P_{0}T^{n}$
 \end{prop}
   
  \begin{demo}
  Par récurrence.
\end{demo}   
   
   \begin{defn}
   
   On appelle \textbf{répartition stable de probabilité} une matrice ligne $P$ dont tous les coefficients sont positifs et de somme 1 vérifiant: 
   
   \begin{center}
   $\boxed{P=PT}$
   \end{center}
   
   \end{defn}
   
   \begin{thm}
   Avec les notations précédentes, si $(p;q)\neq (0;0)$ et $(p;q)\neq (1;1)$ alors:
   
   1) il existe une unique répartition stable de probabilité $P$ donnée par:
   
   \begin{center}
   $\boxed{P=\left( \dfrac{q}{p+q}\hspace{0.4cm}\dfrac{p}{p+q}\right) }$
   \end{center}
   
   2) la suite $(P_{n})$ converge vers $P$, indépendamment   de $P_{0}$
   
   \end{thm}
   
\begin{demo}

   $0\leq p \leq 1$ et $0\leq q \leq 1$ donc: $0\leq p+q\leq 2$
   
   Comme $p+q\neq 0$ et $p+q\neq 2$ on a: $0 < p+q<  2$ puis $-2<-p-q<0$ et enfin:
   
   \begin{center}
   $\boxed{-1<1-p-q<1}$ (*)
   \end{center}
   
   On a vu que: $a_{n+1}=(1-p)a_{n}+qb_{n}$ et $b_{n}=1-a_{n}$ donc: $a_{n+1}=(1-p-q)a_{n}+q$
   
   A partir de cette dernière formule, on peut démontrer par récurrence que:
   
   \begin{center}
   $\boxed{a_{n}=(1-p-q)^{n}\left( a_{0}-\dfrac{q}{p+q}\right) +\dfrac{q}{p+q}}$ puis $\boxed{b_{n}=\dfrac{p}{p+q}-(1-p-q)^{n}\left( a_{0}-\dfrac{q}{p+q}\right) }$
   \end{center}
   
   D'après (*), $\lim_{n \to +\infty}\limits a_{n}=\dfrac{q}{p+q}$ et $\lim_{n \to +\infty}\limits b_{n}=\dfrac{p}{p+q}$. On a bien: $\lim_{n \to +\infty}\limits P_{n}=P$
   
   \begin{center}
    \begin{tabular}{cc}
    &$ \begin{pmatrix}
    \hspace{1.1cm}1-p & p \textcolor{white}{.} \\ 
  \hspace{1.1cm} q & 1-q \textcolor{white}{.} 
   \end{pmatrix} $ \\ 
 $PT=P\Leftrightarrow=\begin{pmatrix}
     x & y
   \end{pmatrix} $ & $\begin{pmatrix}
    (1-p)x+qy & px+(1-q)y  
   
   \end{pmatrix}= (x \hspace{0.3cm}y) $ \\
   
 \end{tabular} 
   \end{center}
     
   On obtient le système: $ \left\lbrace
\begin{array}{l}

 x-xp+yq = x\\[0.3cm]
 
  x+y = 1\\[0.3cm] 
\end{array}
\right.  $ qui donne: $ \left\lbrace
\begin{array}{l}

 x=\dfrac{q}{p+q}\\[0.3cm]
 
  y=\dfrac{p}{p+q}\\[0.3cm] 
\end{array}
\right.  $
   
      
   \end{demo}   
      
\section{Marche aléatoire entre plusieurs états}
 
  Les définitions et propriétés précédentes se généralisent à un système qui peut se trouver dans plusieurs états.
  
\begin{thm}
 
 Si la matrice de transition $T$ a une puissance n'ayant aucun coefficient nul, alors:
 
 1) il existe une unique répartition stable de probabilité $P$ telle que $PT=P$
   
      
   2) la suite $(P_{n})$ converge vers $P$, indépendamment   de $P_{0}$ \end{thm}


\chapter{Les graphes}
\section{Premi\`eres notions}

De très nombreux problèmes pratiques peuvent être ainsi schématisés à l'aide d'un graphe ; en simplifiant la représentation, on peut ainsi trouver plus rapidement la solution (ou en voir
l'impossibilité !). Pour certains problèmes, comme ceux que nous venons de voir, les arêtes n'ont pas d'orientation. Pour d'autres, il est indispensable d'avoir une orientation sur le graphe : le plan d'une ville comme graphe non orienté satisfera le piéton, tandis que ce même graphe orienté par les sens de circulation sera bien plus apprécié de l'automobiliste.

Une question importante est celle du choix du graphe associé à une situation donnée (il peut
y en avoir plusieurs) ; comment choisir les sommets et les arêtes ? Comme on vient de le voir
dans le problème des segments, ce n'est pas toujours évident. Dans des paragraphes ultérieurs, on étudiera des questions de compatibilité, il faudra décider si les arêtes correspondent aux couples de points compatibles ou incompatibles, et si les arêtes sont orientées ou non.

Nous allons formaliser les notions qui précèdent.

%%%%%%%%%%%%%%%%%%%%%%%%%%%%%%%%%%%%%%%%%%%%%%%%%%


\begin{defn}[Graphe, sommets, arêtes, sommets adjacents] Un graphe $G$ (non orienté) est constitué d'un ensemble $S$ de points appelés \emph{sommets}, et d'un ensemble $A$ d'\emph{arêtes}, tels qu'à chaque arête sont associés deux sommets, appelés ses \emph{extrémités}.\\
Deux sommets qui sont les extrémités d'une arête sont dits \emph{adjacents}. \end{defn}


Les deux extrémités peuvent être distinctes ou confondues ; dans ce dernier cas, l'arête s'appelle une \emph{boucle}. Deux arêtes peuvent aussi avoir les mêmes extrémités (on dit alors qu'elles sont \emph{parallèles}). Cependant, la très grande majorité des problèmes que nous rencontrerons, où des graphes non orientés seront en jeu, concerne des graphes \emph{simples}, c'est-à-dire sans boucles ni arêtes parallèles. Les termes \emph{simples} et \emph{parallèles} ne sont pas à retenir.


\begin{exemple}  On considère le graphe $G_1$, de la figure \ref{lexemple}.
Le sommet $s_4$ est un sommet isolé, l'arête $a$ est une boucle, $b$ et $c$ sont des arêtes ayant mêmes extrémités, les sommets $s_1$ et $s_2$ sont adjacents, ainsi que $s_1$ et $s_3$,
puisqu'ils sont reliés par une arête.\\

\centering
%\label{lexemple}\caption{Le graphe $G_1$}
\psset{xunit=1cm , yunit=1cm}
\begin{pspicture*}(-0.7,-0.7)(5.2,3.7)
\def\xmin{-0.5} \def\xmax{5} \def\ymin{-0.5} \def\ymax{3.5}

\psset{linecolor=black, linewidth=.5pt, arrowsize=2pt 4}
\begin{psmatrix}[mnode=circle]
$s_1$ &  & $s_2$ \\
$s_3$ &  & $s_4$
\end{psmatrix}
\psset{shortput=nab}
\nccircle[angleA=0]{1,1}{0.5cm}^{$a$}
\ncarc[arcangle=15]{1,1}{1,3}^{$b$}
\ncarc[arcangle=15]{1,3}{1,1}^{$c$}
\ncline{1,1}{2,1}^{$d$}
\ncline{1,3}{2,1}^{$e$}
\end{pspicture*}

\end{exemple}


\begin{rem}
 La position des sommets et la longueur ou l'allure des arêtes n'ont aucune importance. 
\end{rem}



\begin{exo}
Parmi les sept graphes donnés dans la figure, déterminer ceux qui sont identiques.
%\begin{figure}[!t]
\centering
%\caption{Sept graphes}\label{septgraphes}
\psset{xunit=1.15cm , yunit=1.15cm}
\begin{pspicture*}(-1.2,-1.2)(12.7,4.7)
\def\xmin{-1} \def\xmax{12.5} \def\ymin{-1} \def\ymax{4.5}
\psset{linecolor=black, linewidth=.5pt, arrowsize=2pt 4}
\psline{o-o}(1.0000,0.0000)(2.0000,2.0000)
\psline{o-o}(1.0000,0.0000)(1.0000,4.0000)
\psline{o-o}(1.0000,0.0000)(0.0000,2.0000)
\psline{o-o}(2.0000,2.0000)(1.0000,4.0000)
\psline{o-o}(2.0000,2.0000)(0.0000,2.0000)
\psline{o-o}(1.0000,4.0000)(0.0000,2.0000)
\psline{o-o}(3.0000,1.0000)(5.0000,1.0000)
\psline{o-o}(3.0000,1.0000)(4.0000,2.0000)
\psline{o-o}(3.0000,1.0000)(4.0000,4.0000)
\psline{o-o}(5.0000,1.0000)(4.0000,2.0000)
\psline{o-o}(5.0000,1.0000)(4.0000,4.0000)
\psline{o-o}(4.0000,2.0000)(4.0000,4.0000)
\psline{o-o}(6.0000,2.0000)(8.0000,1.0000)
\psline{o-o}(8.0000,1.0000)(9.0000,3.0000)
\psline{o-o}(9.0000,3.0000)(7.0000,4.0000)
\psline{o-o}(7.0000,4.0000)(6.0000,3.0000)
\psline{o-o}(6.0000,3.0000)(6.0000,2.0000)
\psline{o-o}(10.0000,0.0000)(12.0000,3.0000)
\psline{o-o}(10.0000,0.0000)(11.0000,4.0000)
\psline{o-o}(12.0000,0.0000)(10.0000,3.0000)
\psline{o-o}(12.0000,0.0000)(11.0000,4.0000)
\psline{o-o}(12.0000,3.0000)(10.0000,3.0000)
\rput(1,-0.5){$G_1$}
\rput(4,-0.5){$G_2$}
\rput(7.5,-0.5){$G_3$}
\rput(11,-0.5){$G_4$}
\end{pspicture*}

\medskip

\psset{xunit=1.15cm , yunit=1.15cm}
\begin{pspicture*}(-1.2,-1.2)(12.7,4.7)
\def\xmin{-1} \def\xmax{12.5} \def\ymin{-1} \def\ymax{4.5}
\psset{linecolor=black, linewidth=.5pt, arrowsize=2pt 4}
\psline{o-o}(1.0000,1.0000)(3.0000,1.0000)
\psline{o-o}(1.0000,1.0000)(2.0000,2.0000)
\psline{o-o}(1.0000,1.0000)(1.0000,4.0000)
\psline{o-o}(3.0000,1.0000)(2.0000,2.0000)
\psline{o-o}(3.0000,1.0000)(3.0000,4.0000)
\psline{o-o}(2.0000,2.0000)(2.0000,3.0000)
\psline{o-o}(2.0000,3.0000)(3.0000,4.0000)
\psline{o-o}(2.0000,3.0000)(1.0000,4.0000)
\psline{o-o}(3.0000,4.0000)(1.0000,4.0000)
\psline{o-o}(5.0000,1.0000)(6.0000,0.0000)
\psline{o-o}(5.0000,1.0000)(7.0000,3.0000)
\psline{o-o}(5.0000,1.0000)(5.0000,3.0000)
\psline{o-o}(6.0000,0.0000)(7.0000,1.0000)
\psline{o-o}(6.0000,0.0000)(6.0000,4.0000)
\psline{o-o}(7.0000,1.0000)(7.0000,3.0000)
\psline{o-o}(7.0000,1.0000)(5.0000,3.0000)
\psline{o-o}(7.0000,3.0000)(6.0000,4.0000)
\psline{o-o}(6.0000,4.0000)(5.0000,3.0000)
\psline{o-o}(10.0000,0.0000)(11.0000,1.0000)
\psline{o-o}(10.0000,0.0000)(10.0000,4.0000)
\psline{o-o}(10.0000,0.0000)(9.0000,1.0000)
\psline{o-o}(11.0000,1.0000)(11.0000,3.0000)
\psline{o-o}(11.0000,1.0000)(9.0000,1.0000)
\psline{o-o}(11.0000,3.0000)(10.0000,4.0000)
\psline{o-o}(11.0000,3.0000)(9.0000,3.0000)
\psline{o-o}(10.0000,4.0000)(9.0000,3.0000)
\psline{o-o}(9.0000,3.0000)(9.0000,1.0000)
\rput(2,-0.5){$G_5$}
\rput(6,-0.5){$G_6$}
\rput(10,-0.5){$G_7$}
\end{pspicture*}
%\end{figure}
\end{exo}



\begin{rem} C'est un problème très difficile en général, dès que le nombre de sommets est assez
grand.
\end{rem}

Une première manière d'évaluer la complication d'un graphe est de compter le nombre de ses
sommets :

\begin{defn}[Ordre d'un graphe] L'\emph{ordre} d'un graphe est le nombre de ses sommets.\end{defn}

\begin{defn}[Degré d'un sommet, parité d'un sommet] On appelle \emph{degré d'un sommet} le nombre d'arêtes dont ce sommet est une extrémité (les boucles étant comptées deux fois). Un sommet est \emph{pair} (respectivement \emph{impair}) si son degré est un nombre pair (respectivement impair). \end{defn}

\begin{exemple} Dans la figure, $s_1$ est de degré 5, $s_2$ de degré 3, $s_4$ de degré 0. \end{exemple}


On prouve facilement le théorème suivant :

\begin{thm} La somme des degrés de tous les sommets d'un graphe est égale à deux fois le nombre d'arêtes de ce graphe ; c'est donc un nombre pair. \end{thm}

\begin{demo} Lorsque on additionne les degrés des sommets, chaque arête est comptée deux fois, une fois pour chaque extrémité. \end{demo}

\begin{prop} Dans un graphe le nombre de sommets impairs est toujours pair. \end{prop}

\begin{demo} En effet, sinon, la somme des degrés des sommets serait impaire. \end{demo}

\begin{exo} À l'aide de ce théorème ou de cette propriété, montrer que certains des problèmes donnés en introduction n'ont pas de solution. \end{exo}

%\sautpage

\section{Graphes complets}

\begin{defn}[Graphe complet] Un graphe (simple) est dit \emph{complet} si tous ses sommets sont \emph{adjacents}, c'est-à-dire si toutes les arêtes possibles existent. On appelera $K_n$ le graphe complet à $n$ sommet (il n'y en a qu'un). \end{defn}

\begin{exo*}
Parmi les graphes de la figure , lesquels sont complets ?
\end{exo*}

\begin{exo*}
Quel est le degré de chacun des sommets d'un graphe complet d'ordre $n$ ?
\end{exo*}

\section{Sous-graphes}

\begin{defn}[Sous-graphe, sous-graphe engendré par des sommets]
Soit $G$ un graphe, le graphe $G'$ est \underline{un} \emph{sous-graphe} de $G$, si :
\begin{itemize}
	\item l'ensemble des sommets de $G'$ est inclu dans celui des sommets de $G$ ;
	\item l'ensemble des arêtes de $G'$ est inclu dans celui des arêtes de $G$.
\end{itemize}
Si, de plus, les arêtes de $G'$ sont exactement \emph{\underline{toutes}} les arêtes de $G$ joignant les sommets de $G'$, on dit que $G'$ est \emph{\underline{le} sous-graphe de $G$ engendré par $s_1, s_2, \dots , s_n$}.
\end{defn}

\begin{exemple} Considérons le graphe $G$ dont les sommets sont les villes françaises possédant une gare et dont les arêtes sont les voies ferrées reliant ces villes (on excluera les gares où ne passent plus de voies).
\begin{itemize}
	\item Le graphe $G'$ dont les sommets sont les villes d'un même département possédant une gare et dont les arêtes sont les voies ferrées reliant ces villes est le sous-graphe de $G$ engendré par ces villes.
	\item Le graphe $G''$ dont les sommets sont les villes où passe un TGV et dont les arêtes sont des voies ferrées TGV est un simple sous-graphe de $G$ car il existe des voies normales reliant, par exemple, Marseille à Lyon, qui ne sont pas dans ce sous-graphe.
\end{itemize}
\end{exemple}

\begin{exemple}
La figure présente deux sous graphes du graphe $G_1$.



\begin{multicols}{2}

\begin{center}
\psset{xunit=1cm , yunit=1cm}
\begin{pspicture*}(-0.7,-0.7)(5.2,3.7)
\psset{linecolor=black, linewidth=.5pt, arrowsize=2pt 4}
\begin{psmatrix}[mnode=circle]
$s_1$ &  & $s_2$ \\
$s_3$ &  &
\end{psmatrix}
\psset{shortput=nab}
\nccircle[angleA=0]{1,1}{0.5cm}^{$a$}
\ncarc[arcangle=15]{1,3}{1,1}^{$c$}
\end{pspicture*}

$G_2$ un sous-graphe de $G_1$

\end{center}

\begin{center}
\psset{xunit=1cm , yunit=1cm}
\begin{pspicture*}(-0.7,-0.7)(5.2,3.7)
\def\xmin{-0.5} \def\xmax{5} \def\ymin{-0.5} \def\ymax{3.5}
\psset{linecolor=black, linewidth=.5pt, arrowsize=2pt 4}
\begin{psmatrix}[mnode=circle]
$s_1$ &  & $s_2$ \\
$s_3$ &  &
\end{psmatrix}
\psset{shortput=nab}
\nccircle[angleA=0]{1,1}{0.5cm}^{$a$}
\ncarc[arcangle=15]{1,1}{1,3}^{$b$}
\ncarc[arcangle=15]{1,3}{1,1}^{$c$}
\ncline{1,1}{2,1}^{$d$}
\ncline{1,3}{2,1}^{$e$}
\end{pspicture*}
$G_3$ le sous-graphe de $G_1$ \\ engendré par $s_1, s_2, s_3$
\end{center}
\end{multicols}

\end{exemple}

\begin{exo*}
Dessiner les graphes suivants :\\
\vspace{-1em}\begin{multicols}{2}\begin{itemize}
	\item $G_4$ le sous-graphe de $G_1$ engendré par $s_1, s_2$ ;
	\item $G_5$ le sous-graphe de $G_1$ engendré par $s_1, s_3, s_4$ ;
	\item $G_6$ le sous-graphe de $G_1$ engendré par $s_1, s_4$ ;
	\item $G_7$ le sous-graphe de $G_1$ engendré par $s_2$ et $s_4$.
\end{itemize}\end{multicols}
\end{exo*}

Enfin on a :

\begin{defn}[Sous-graphe stable] On dit qu'un sous-ensemble de l'ensemble des sommets est \emph{stable} s'il ne contient pas de paire de sommets adjacents.
On peut aussi parler de sous-graphe stable : cela revient au même, puisque si un ensemble de
sommets est stable, le graphe engendré, par définition, n'a pas d'arête. \end{defn}

\begin{exemple} Dans l'exercice ci-dessus, le sous-graphe $G_7$ est stable. \end{exemple}

Ce terme de \og stable \fg{} peut paraître arbitraire. Il est en fait naturel si l'on considère ce qu'on appelle un \og graphe d'incompatibilité \fg{} : dans un groupe d'individus, on peut définir un graphe en reliant par une arête les individus qui ne peuvent se supporter. Si l'on veut choisir un sous-groupe de personnes qui travaillent ensemble, il est préférable de choisir un sous-ensemble stable ! On verra en particulier beaucoup d'applications de cette notion dans le paragraphe sur les colorations.

\section{Chaînes et connexité}

Dans bien des problèmes de graphes, il est naturel de considérer ce que l'on peut appeler, de façon informelle, des \og parcours \fg{} ou \og chemins \fg. Le mot utilisé en théorie des graphes est \emph{chaîne}.

La notion intuitive de chaîne, ou plus tard de chaîne orientée, se comprend bien sur un dessin, il est moins facile d'en donner une définition effective.

\begin{defn}[Chaîne, longueur d'une chaîne, cycle]
Une \emph{chaîne} dans un graphe $G$ est une suite finie : $s_0; a_1; s_1; a_2; s_2; a_3; s_3; \dots a_n; s_n$ débutant et finissant par un sommet, alternant sommets et arêtes de telle manière que chaque arête soit encadrée par ses sommets extrémités. \\
La \emph{longueur} de la chaîne est égale au nombre d'arêtes qui la constituent ; la chaîne est \emph{fermée} si $s_0 = s_n$, si de plus toutes ses arêtes sont distinctes on dit alors que c'est un \emph{cycle}.
\end{defn}

\begin{rem} Quand il n'y a pas d'ambiguïté (pas d'arêtes multiples), on peut définir une chaîne
par seulement la suite de ses sommets ou par seulement la suite de ses arêtes. \end{rem}

Enfin on a :

\begin{defn}[Graphe connexe] Un graphe est \emph{connexe} si deux sommets quelconques sont reliés par une chaîne. \end{defn}

\begin{defn}[Distance entre deux sommets] Soit $G$ un graphe connexe, $s$ et $s'$, deux sommets quelconques de $G$. Le graphe étant connexe, il existe au moins une chaîne reliant $s$ et $s'$. On appelle \emph{distance entre $s$ et $s'$} la plus petite des longueurs des chaînes reliant $s$ à $s'$. \end{defn}

\begin{rem} Lorsque le graphe n'est pas connexe, il existe au moins deux sommets qui ne sont pas reliés par une chaîne. On dit parfois que la distance entre ces sommets est infinie.
\end{rem}

\begin{defn}[Diamètre d'un graphe] On appelle \emph{diamètre} d'un graphe connexe, la plus grande distance entre ses sommets. \end{defn}

\begin{rem} Lorsque le graphe n'est pas connexe, on dit parfois que son diamètre est infini.
\end{rem}

%%%%%%%%%%%%%%%%%%%%%%%%%%%%%%%%%%%%%%%%%%%%%%%%%%%%%%%%%%%

\section{Graphes orientés}

Il existe de nombreux domaines où les graphes sont orientés. Par exemple : plan de ville, avec les
sens interdits ; parcours en montagne, où il est utile d'indiquer le sens de montée ! Circuit électrique
en courant continu, où il faut orienter les arêtes pour décider du signe de l'intensité : ce n'est pas
la même chose de faire passer 10 ampères de A vers B ou de B vers A ; graphe d'ordonnancement,
où les arêtes relient une tâche à une autre qui doit la suivre : on ne peut faire la peinture avant le
plâtre.

\begin{defn}
On appelle graphe orienté un graphe où chaque arête est orientée, c'est-à-dire
qu'elle va de l'une des ses extrémités, appelée \emph{origine} ou \emph{extrémité initiale} à l'autre, appelée
\emph{extrémité terminale}.
\end{defn}



Dans un graphe orienté, chaque arête orientée possède un début et une fin.
Toutes les notions que nous avons définies pour un graphe ont un équivalent pour un graphe
orienté. Nous nous contenterons de rajouter le mot \og orienté \fg{} pour préciser ; le contexte rendra
évidente l'interprétation à donner.

En particulier, une chaîne orientée est une suite d'arêtes telle que l'extrémité finale de chacune
soit l'extrémité initiale de la suivante. On prendra garde au fait que l'on peut définir et utiliser
des chaînes (non orientées) sur un graphe orienté. Par exemple, sur un plan de ville où toutes les
rues sont en sens unique, un parcours de voiture correspond à une chaîne orientée, un parcours de
piéton correspond à une chaîne (non orientée).


\end{document}

